\documentclass[14pt, a5paper, twoside, extrafontsizes]{memoir}

\def \sederTitle{משניות סדר נשים}
\def \tractateTitle{סדר נשים}
\usepackage{lettrine}

% indentation of the first word on the first line after the drop capital
\setlength{\DefaultFindent}{6pt}

% extra indentation for the first word on the next line
\setlength{\DefaultNindent}{0pt}

\renewcommand*{\DefaultLoversize}{-.2}
\renewcommand*{\DefaultLraise}{.6}
\renewcommand{\LettrineTextFont}{\relax}

\renewcommand*{\LettrineFont}{%
 \LettrineFontHook
 \fontsize{18pt}{18pt}%
 \selectfont
}

% \leadword{first word [to be enlarged]}{second word}
\newcommand{\leadword}[2]{%
  \lettrine[realheight=true]{#1}{#2}
}


\usepackage{graphicx}
\usepackage{pgffor}
\usepackage{anysize}
\usepackage[all]{nowidow}
\usepackage{relsize}
\usepackage{ragged2e}

\usepackage[pdfauthor={יהודה הנשיא},
  pdftitle=\sederTitle,
  pdfsubject={תורה שבעל פה},
  pdfproducer={סיימון מונטגיו, הוצאת חטף סגול},
  pdfdirection=R2L,
  pdfcreator=xelatex,
  colorlinks=false, linkcolor=blue,
  linkbordercolor={.6 .6 .6},
  plainpages=true]{hyperref}
\makeatletter
\Hy@AtBeginDocument{%
  % Override border definition set with colorlinks=true
  \def\@pdfborder{0 0 1}
  % Overrides border style set with colorlinks=true
  % Hyperlink border style will be underline of width 1pt
  \def\@pdfborderstyle{/S/U/W 1}
}
\makeatother

\usepackage{polyglossia}
\setmainlanguage[numerals=hebrew]{hebrew}
\setmainfont{SBL Hebrew}
\setotherlanguage{english}

% General formatting
\marginsize{1.2in}{.8in}{.5in}{.5in}
\renewcommand{\baselinestretch}{1.1}
\parskip 8pt
\parindent 0pt
\makeevenhead{myheadings}{\thepage}{}{\texthebrew{\sederTitle}}
\makeoddhead{myheadings}{\texthebrew{\tractateTitle}}{}{\thepage}
\makepagestyle{mytitlepage}

\newcounter{perekCount}
\newcounter{halachaCount}[perekCount]
\newcounter{tractateCount}
\newcounter{perakim}

% \tractatehead{title for printing}{hyperlink reference}
\newcommand{\tractatehead}[2]{%
 \cleardoublepage
 \setcounter{perekCount}{0}
 \thispagestyle{empty}
 \def \tractateTitle{#1}
{\centering\LARGE\bfseries \hypertarget{#2}{#1}

  }}

% \summarycontentsline{tractate name}{hyperlink reference}
%                     {number of chapters [numeric]}
\newcommand{\summarycontentsline}[3]{%
\setcounter{perakim}{#3}%
\stepcounter{tractateCount}%
\textlarger{\alph{tractateCount}} \hyperlink{#2}{#1} \alph{perakim} פרקים\hdot%
}

\newcommand{\summarycontents}[1]{
  \centerlastline
  \foreach \linkText / \linkName / \numChapters in #1 {
      \summarycontentsline{\linkText}{\linkName}{\numChapters}
  } 
}

% \sederContents{number of tractates [hebrew]}
%               {list of tractates}{total number of chapters [hebrew]}
\newcommand{\sederContents}[3]{
  {
    \centering
    {\Large משניות \sederTitle}

  יש בו #1 מסכתות וזהו סדורן

    \begin{minipage}{0.9\textwidth}
    \small
    \summarycontents{#2}

    בין הכל \textlarger[2]{#3} פרקים׃
    \end{minipage}

  }
}

% centered rubric
\newcommand{\crubric}[1]{{\centering\small\bfseries{#1}

  }}

% centered header (keep with next paragraph)
\newcommand{\chead}[1]{%
  \vspace{0.25em}
  \needspace{4\baselineskip}
  \crubric{#1}
}

% centered footer
\newcommand{\cfoot}[1]{%
  \vspace{0.25em}
  \crubric{#1}
}

\newcommand{\mispar}[1]{{\footnotesize\textsuperscript{\alph{#1}}}}

% There is no Unicode character for the raised dot used in many Hebrew books
% as mid-sentence punctuation. This is a hack to replace it using
% U+005C4 HEBREW MARK UPPER DOT
\sfcode`\^^^^05c43000
\newcommand{\hdot}{\raisebox{-0.5\height}{^^^^2009^^^^2009^^^^05c4} }

\tolerance=8192

% \perek{first word}{rest of first halacha}
% First halacha in a perek
\newcommand{\perek}[2]{%
 \stepcounter{perekCount}
 \perekbar{פרק \alph{perekCount}}{#1}{#2}
}

% \baraita[subhead]{first word}{rest of first halacha}
% First halacha in a baraita (e.g. after the last perek of Bikkurim)
\newcommand{\baraita}[3][]{%
 \setcounter{halachaCount}{0}
 \perekbar[#1]{ברייתא}{#2}{#3}
}

% \perekbar[subhead]{heading}{first word}{rest of first halacha}
% Called by \perek and \baraita
\newcommand{\perekbar}[4][]{%
 \vspace{0.5em}
 \stepcounter{halachaCount}
 \needspace{2\baselineskip}
 \crubric{#2}

 \centerlastline
 #1

 \leadword{#3}#4
}

% \halacha{text}
% Prints a halacha
\newcommand{\halacha}[1]{%
 \stepcounter{halachaCount}
 \mispar{halachaCount}
 \centerlastline
 #1\nowidow
}

\newcommand{\selik}{~~{\footnotesize סליק}}

% \titlepage{Hebrew chronogram for year of issue}
%           {Gregorian year of issue}
\newcommand{\titlepage}[2]{%
\Centering

  {\huge\bfseries משניות}

  {\LARGE\bfseries\sederTitle}

  \vspace{1em}

  \vspace{0.5em}
  על פי מהדורת הרב דוד אלטאראס ז״ל\\שהודפס בשנת ה׳תק״ע בעיר פיסא

    \vspace{8em}
    פה {\LARGE ירושלים} יע״א

    {\small שנת} #1 {\small לפ״ק}

    \vspace{2em}
      \includegraphics[width=10mm]{hatafSegolLogoNoText.png}\\

    \vspace{.5em}
  {
    \bfseries הוצאת חטף־סגול\\#2

 {\footnotesize  הובא לדפוס באמצעות \XeLaTeX\ ע״י שמעון מונטגיו}

  }

}

% code for special handling of specific numbers by "Circumscribe" copied from https://tex.stackexchange.com/questions/300008/modify-specific-hebrew-alpha-numerals-on-page-number
% with a few modifications to the function names
\makeatletter
\let\@hebrew@numeral@usual\@hebrew@numeral  %% <- store old definition
\newcommand*\@hebrew@numeral@special[1]{%   %% <- new definition
  \ifcsdef{specialnum@\number#1}            %% <- if this number is listed
    {\csuse{specialnum@\number#1}}          %% <- give it special handling
    {\@hebrew@numeral@usual{#1}}%           %% <- otherwise use usual handling
}
\let\@hebrew@numeral\@hebrew@numeral@special %% <- replace old definition
\makeatother

%% Declare numbers for special handling:
\newcommand*\newspecialnum[2]{\csdef{specialnum@#1}{#2}}
\newspecialnum{16}{יו}    % At least some Livorno printings don't use טז for 16
\newspecialnum{116}{קיו}
\newspecialnum{216}{ריו}
\newspecialnum{316}{שיו}
\newspecialnum{416}{תיו}
\newspecialnum{516}{תקיו}
\newspecialnum{616}{תריו}
\newspecialnum{716}{תשיו}
\newspecialnum{816}{תתיו}
\newspecialnum{916}{תתקיו}
\newspecialnum{18}{חי}    % Add some other "specials" while we're here
\newspecialnum{270}{ער}
\newspecialnum{272}{ערב}
\newspecialnum{275}{ערה}



\begin{document}
\frontmatter
\pagestyle{myheadings}
\thispagestyle{mytitlepage}
\begin{hebrew}
  \pagenumbering{arabic}
{\Centering

  {\huge\bfseries משניות}

  {\LARGE\bfseries סדר נשים}

  \vspace{1em}

  \vspace{0.5em}
  על פי מהדורת הרב דוד אלטאראס ז״ל\\שהודפס בשנת ה׳תק״ע בעיר פיסא

    \vspace{8em}
    פה {\LARGE ירושלים} יע״א

    {\small שנת}
    {\Large שׄלׄוׄםׄ}\ בין {\Large אׄ}יש ל{\Large אׄ}ש{\Large תׄוׄ}
    {\small לפ״ק}

    \vspace{2em}
  \includegraphics[width=10mm]{hatafSegolLogoNoText.png}\\

    \vspace{.5em}
  {
    \bfseries הוצאת חטף־סגול\\2024

 {\footnotesize  הובא לדפוס באמצעות \XeLaTeX\ ע״י שמעון מונטגיו}

  }

}

\clearpage
\thispagestyle{empty}

\centerlastline
{
\small
\parskip 2pt
{\centering
  \scriptsize
  תפלת ר׳ נחוניא בן הקנה שחייב לאומרה כל הנכנס לבית המדרש ללמוד

}

יְהִי רָצוֹן מִלְּפָנֶיךָ יְיָ אֱלֹהַי שֶׁלֹּא יֶאֱרַע דְּבַר תְּקָלָה עַל יָדִי וְלֹא אֶכָּשֵׁל בִּדְבַר הֲלָכָה
וְיִשְׂמְחוּ בִי חֲבֵירַי וְלֹא אוֹמַר עַל טָהוֹר טָמֵא וְלֹא עַל טָמֵא טָהוֹר וְלֹא עַל מֻתָּר אָסוּר
וְלֹא עַל אָסוּר מֻתָּר וְלֹא יִכָּֽשְׁלוּ חֲבֵירַי בִּדְבַר הֲלָכָה וְאֶשְׂמַח בָּהֶם׃

{\centering
  {\scriptsize
    והאר״י זצוק״ל היה אומר אחר תפלה זו שני פסוקים אלו

    }
    כִּי יְיָ יִתֵּן חָכְמָה מִפִּיו דַּעַת וּתְבוּנָה׃\\
    גַּל עֵינַי וְאַבִּיטָה נִפְלָאוֹת מִתּוֹרָתֶךָ׃

    \vspace{1em}
   {\larger קודם הלימוד}

    {\scriptsize כתוב בספר תוצאות חיים כשפותח הספר ללמוד בו יאמר}
}

\begin{tp}{396}
הִנְּנִי רוֹצֶה לִלְמוֹד כְּדֵי שֶׁיְּבִיאֵנִי תַּלְמוּד זֶה לִידֵי מַעֲשֶׂה וְלִידֵי מִדּוֹת יְשָׁרוֹת\hdot
וְלִידֵי יְדִיעַת הַתּוֹרָה\hdot
וַהֲרֵינִי עוֹשֶׂה לְשֵׁם יִחוּד קֻדְשָׁא בְרִיךְ הוּא וּשְׁכִינְתֵּיהּ בִּדְחִילוּ וּרְחִימוּ\hdot
וִיהִי נוֹעַם יְהֹוָה אֱלֹהֵינוּ עָלֵינוּ וּמַעֲשֵׂה יָדֵינוּ כּוֹנְנָה עָלֵינוּ וּמַעֲשֵׂה יָדֵינוּ כּוֹנְנֵהוּ׃
\end{tp}

}

\vspace{1em}
{
\parskip 2pt
\sederContents{נשים}{ז׳}{{יבמות/yebamot/16,
  כתובות/ketubot/13,
  קדושין//12,
  גיטין//9,
  סוטה/sota/9,
  נדרים/nedarim/11,
  נזיר/nazir/9}}{ע״א}


\vspace{4pt}
       { \centering
         \small
  {\larger  אבל} כפי סדר תוספת יום טוב כך הוא\\
  יבמות\hdot כתובות\hdot נדרים\hdot נזיר\hdot סוטה\hdot גיטין\hdot קדושין\\
  אבל הר״ם סידר גיטין אחר נזיר\\
  ואלו ואלו דברי אלהים חיים׃}

}
\mainmatter
\thispagestyle{empty}

\tractatehead{מסכת יבמות}{yebamot}

\chead{יום יוד}

\perek{חֲמֵשׁ}{עֶשְׂרֵה נָשִׁים פּוֹטְרוֹת צָרוֹתֵיהֶן\hdot
וְצָרוֹת צָרוֹתֵיהֶן\hdot
מִן הַחֲלִיצָה וּמִן הַיִּבּוּם\hdot
עַד סוֹף הָעוֹלָם\hdot
וְאֵלּוּ הֵן\hdot
בִּתּוֹ\hdot
וּבַת בִּתּוֹ\hdot
וּבַת בְּנוֹ\hdot
בַּת אִשְׁתּוֹ\hdot
וּבַת בְּנָהּ\hdot
וּבַת בִּתָּהּ\hdot
חֲמוֹתוֹ\hdot
וְאֵם חֲמוֹתוֹ\hdot
וְאֶם חָמִיו\hdot
אֲחוֹתוֹ מֵאִמּוֹ\hdot
וַאֲחוֹת אִמּוֹ\hdot
וַאֲחוֹת אִשְׁתּוֹ\hdot
וְאֵשֶׁת אָחִיו מֵאִמּוּ\hdot
וְאֵשֶׁת אָחִיו שֶׁלֹּא הָיָה בְעוֹלָמוֹ\hdot
וְכַלָּתוֹ\hdot
הֲרֵי אֵלּוּ פּוֹטְרוֹת צָרוֹתֵיהֶן\hdot
וְצָרוֹת צָרוֹתֵיהֶן מִן הַחֲלִיצָה וּמִן הַיִּבּוּם עַד סוֹף הָעוֹלָם\hdot
וְכֻלָּן אִם מֵתוּ\hdot
אוֹ מִאֲנוּ\hdot
אוֹ נִתְגָּֽרְשׁוּ\hdot
אוֹ שֶׁנִּמְצְאוּ אַיְלוֹנִיּוֹת\hdot
צָרוֹתֵיהֶן מֻתָּרוֹת\hdot
וְאֵי אַתָּה יָכוֹל לוֹמַר בַּחֲמוֹתוֹ (וּבְאֵם חֲמוֹתוֹ) וְבְאֵם חָמִיו\hdot
שֶׁנִּמְצְאוּ אַיּלוֹנִיּוֹת אוֹ שֶׁמִּאֲנוּ׃}

\halacha{כֵּיצִד פּוֹטְרוֹת צָרוֹתֵיהֶן\hdot
הָֽיְתָה בִתּוֹ\hdot
אוֹ אַחַת מִכָּל הָעֲרָיוֹת הָאֵלּוּ\hdot
נְשׂוּאָה לְאָחִיו\hdot
וְלוֹ אִשָּׁה אַחֶרֶת וּמֵת\hdot
כַּשֵּׁם שֶׁבִּתּוֹ פְטוּרָה\hdot
כַּךְ צָרָתָהּ פְּטוּרָה\hdot
הָֽלְכָה צָרַת בִּתּוֹ\hdot
וְנִשֵּׂאת לְאָחִיו הַשֵּׁנִי\hdot
וְלוֹ אִשָּׁה אַחֶרֶת וּמֵת\hdot
כַּשֵּׁם שֶׁצָּרַת בִּתּוֹ פְטוּרָה\hdot
כַּךְ צָרַת צָרָתָהּ פְּטוּרָה\hdot
אֲפִילוּ הֵן מֵאָה\hdot
כֵּיצַד אִם מֵתוּ צָרוֹתֵיהֶן מֻתָּרוֹת\hdot
הָֽיְתָה בִתּוֹ אוֹ אַחַת מִכָּל הָעֲרָיוֹת הָאֵלּוּ\hdot
נְשׂוּאָה לְאָחִיו\hdot
וְלוֹ אִשָּׁה אָחֶרֶת\hdot
מֵתָה בִתּוֹ אוֹ נִתְגָּֽרְשָׁה\hdot
וְאַחַר כַּךְ מֵת אָחִיו\hdot
צָרָתָהּ מֻתֶּרֶת\hdot
וְכָל הַיְּכוֹלָה לְמָאֵן וְלֹא מִיאֲנָה\hdot
צָרָתָהּ חוֹלֶצֶת וְלֹא מִתְיַבֶּמֶת׃}

\halacha{שֵׁשׁ עֲרָיוֹת חֲמוּרוֹת מֵאֵלּוּ\hdot
מִפְּנֵי שֶׁנְּשֹוּאוֹת לַאֲחֵרִים צָרוֹתֵיהֶן מֻתָּרוֹת\hdot
אִמּוֹ\hdot
וְאֵשֶׁת אָבִיו\hdot
וַאֲחוֹת אָבִיו\hdot
אֲחוֹתוֹ מֵאָבִיו\hdot
וְאֵשֶׁת אֲחִי אָבִיו\hdot
וְאֵשֶׁת אָחִיו מֵאָבִיו׃}

\halacha{בֶּית שַׁמָּאי מַתִּירִין הַצָּרוֹת לָאַחִים\hdot
וּבֵית הִלֵּל אוֹסְרִין\hdot
חָֽלְצוּ\hdot
בֵּית שַׁמָּאי פּוֹסְלִין מִן הַכְּהֻנָּה\hdot
וּבֵית הִלֵּל מַכְשִׁירִין\hdot
נִתְיַבְּמוּ בֵּית שַׁמָּאי מַכְשִׁירִין\hdot
וּבֵית הִלֵּל פּוֹסְלִין\hdot
אַף עַל פִּי שֶׁאֵלּוּ (אוֹסְרִין וְאֵלּוּ מַתִּירִין\hdot
אֵלּוּ) פּוֹסְלִין וְאֵלּוּ מַכְשִׁירִין\hdot
לֹא נִמְנְעוּ בֵּית שַׁמָּאי מִלִּשָּׂא נָשִׁים מִבֵּית הִלֵּל\hdot
וְלֹא בֵית הִלֵּל מִבֵּית שַׁמָּאי\hdot
כָּל הַטַּהֲרוֹת וְהַטֻּמְאוֹת שֶׁהָיוּ אֵלּוּ מְטַהֲרִין וְאֵלּוּ מְטַמְּאִין\hdot
לֹא נִמְנְעוּ עוֹשִׂין טַהֲרוֹת\hdot
אֵלּו עַל גַּבֵּי אֵלּוּ׃}

\perek{כֵּיצַד}{אֵשֶׁת אָחִיו שֶׁלֹּא הָיָה בְעוֹלָמוֹ\hdot
שְׁנֵי אַחִים וּמֵת אַחַד מֵהֶם\hdot
וְנוֹלָד לָהֶם אָח\hdot
וְאַחַר כַּךְ יִבֵּם הַשֵּׁנִי אֶת אֵשֶׁת אָחִיו וּמֵת\hdot
הָרִאשׁוֹנָה יוֹצֵאה מִשּׁוּם אֵשֶׁת אָחִיו שֶׁלֹּא הָיָה בְעוֹלָמוֹ\hdot
וְהַשְּׁנִיָּה מִשּׁוּם צָרָתָהּ\hdot
עָשָׂה בָהּ מַאֲמָר וּמֵת\hdot
שְׁנִיָּה חוֹלֶצֶת וְלֹא מִתְיַבֶּמֶת׃}

\halacha{שְֹנֵי אַחִים וּמֵת אַחַד מֵהֶן\hdot
וְיִבֵּם הַשֵּׁנִי אֶת אֵשֶׁת אָחִיו\hdot
וְאַחַר כַּךְ נוֹלָד לָהֶן אָח וּמֵת\hdot
הָרִאשׁוֹנָה יוֹצְאָה מִשּׁוּם אֵשֶׁת אָחִיו שֶׁלֹּא הָיָה בְעוֹלָמוֹ\hdot
וְהַשְּׁנִיָּה מִשּׁוּם צָרָתָהּ\hdot
עָשָׂה בָהּ מַאֲמָר וּמֵת\hdot
הַשְּׁנִיָּה חוֹלֶצֶת וְלֹא מִתְיַבֵֶּמֶת\hdot
רִבִּי שִׁמְעוֹן אוֹמֵר מְיַבֵּם לְאֵיזוֹ מֵהֶן שֶׁיִּרְצֶה\hdot
אוֹ חוֹלֵץ לְאֵיזוֹ מֵהֶן שֶׁיִּרְצֶה׃}

\tractatehead{מסכת כתובות}{ketubot}

\perek{בְּתוּלָה}{נִשֵּׂאת לְיוֹם הָֽרְבִיעִי\hdot
וְאַלְמָנָה לְיוֹם חֲמִישִׁי\hdot
שֶׁפַּעֲמַיִם בְּשַׁבָּת בָּתֵּי דִינִין יוֹשְׁבִין בָּעֲרָיוֹת בַּיּוֹם הַשֵּׁנִי\hdot
וּבַיּוֹם הַחֲמִישִׁי\hdot
שֶׁאִם הָיָה לוֹ טַעֲנַת בְּתוּלִים הָיָה מַשְׁכִּים לְבֵית דִּין׃}

\halacha{בְּתוּלָה כְּתוּבָתָהּ מָאתַיִם\hdot
וְאַלְמָנָה מָנֶה\hdot
בְּתוּלָה אַלְמָנָה גְּרוּשָׁה וַחֲלוּצָה מִן הָאֵרוּסִין\hdot
כְּתוּבָתָן מָאתַיִם וְיֵשׁ לָהֶן טַעֲנַת בְּתוּלִים\hdot
הַגִּיּוֹרֶת וְהַשְּׁבוּיָה וְהַשִּׁפְחָה\hdot
שֶׁנִּפְדּוּ וְשֶׁנִּתְגַּיְּרוּ וְשֶׁנִּשְׁתַּחְרְרוּ\hdot
פְּחוּתוֹת מִבְּנוֹת שָׁלֹשׁ שָׁנִים וְיוֹם אֶחָד\hdot
כּתוּבָתָן מָאתַיִם\hdot
וְיֵשׁ לָהֶן טַעֲנַת בּתוּלִים׃}

\halacha{הַגָּדוֹל שֶׁבָּא עַל הַקְּטָנָה\hdot
וְהַקָּטָן שֶׁבָּא עַל הַגְּדוֹלָה\hdot
וּמֻכַּת עֵץ\hdot
כְּתוּבָתָן מָאתַיִם\hdot
דִּבְרֵי רִבִּי מֵאִיר\hdot
וַחֲכָמִים אוֹמְרִים מֻכַּת עֵץ כְּתוּבָתָהּ מָנֶה׃}

\halacha{בְּתוּלָה אַלְמָנָה גְּרוּשָׁה וַחֲלוּצָה מִן הַנִּשּׂוּאִין\hdot
כְּתוּבָתָן מָנֶה\hdot
וְאֵין לָהֶן טַעֲנַת בְּתוּלִים\hdot
הַגִּיּוֹרֶת וְהַשְּׁבוּיָה וְהַשִּׁפְחָה\hdot
שֶׁנִּפְדּוּ וְשֶׁנִּתְגַּיְּרוּ וְשֶׁנִּשְׁתַּחְרְרוּ\hdot
יְתֵירוֹת עַל בְּנוֹת שָׁלֹשׁ שָׁנִים וְיוֹם אֶחָד\hdot
כְּתוּבָתָן מָנֶה וְאֵין לָהֶן טַעֲנַת בְּתוּלִים׃}

\halacha{הָאוֹכֵל אֵצֶל חָמִיו בִּיהוּדָה שֶׁלֹּא בְעֵדִים\hdot
אֵינוֹ יָכוֹל לִטְעוֹן טַעֲנַת בְּתוּלִים מִפְּנֵי שֶׁמִּתְיַחֵד עִמָּהּ\hdot
אַחַת אַלְמָנַת יִשְׂרָאֵל וְאַחַת אַלְמָנַת כֹּהֵן כְּתוּבָתָהּ מָנֶה\hdot
בֵּית דִּין שֶׁל כֹּהֲנִים גּוֹבִין לַבְּתוּלָה אַרְבַּע מֵאוֹת זוּז\hdot
וְלֹא מִיחוּ בְיָדָם חֲכָמִים׃}

\halacha{הַנּוֹשֵׂא אֶת הָאִשָּׁה וְלֹא מָצָא לָהּ בְּתוּלִים\hdot
הִיא אוֹמֶרֶת מִשֶּׁאֵרַסְתָּנִי נֶאֱנַסְתִּי וְנִסְתַּחֲפָה שָׁדֶךָ\hdot
וְהַלָּה אוֹמֵר לֹא כִי אֶלָּא עַד שֶׁלֹּא אֵרַסְתִּיךְ\hdot
וְהָיָה מִקְחִי מֶקַח טָעוּת\hdot
רַבָּן גַּמְלִיאֵל וְרִבִּי אֱלִיעֶזֶר אוֹמְרִים נֶאֱמֶנֶת\hdot
רִבִּי יְהוֹשֻׁעַ אוֹמֵר לֹא מִפִּיהָ אָנוּ חַיִּין\hdot
אֶלָּא הֲרֵי זוֹ בְּחֶזְקַת בְּעוּלָה עַד שֶׁלֹּא תִתְאָרֵס וְהִטְעַתּוּ עַד שֶׁתָּבִיא רְאָיָה לִדְבָרֶיהָ׃}

\halacha{הִיא אוֹמֶרֶת מֻכַּת עֵץ אָנִי\hdot
וְהוּא אוֹמֵר לֹא כִי\hdot
אֶלָּא דְּרוּסַת אִישׁ אַתְּ\hdot
רַבָּן גַּמְלִיאֵל וְרִבִּי אֱלִיעֶזֶר אוֹמְרִים נֶאֱמֶנֶת\hdot
רִבִּי יְהוֹשֻׁעַ אוֹמֵר לֹא מִפִּיהָ אָנוּ חַיִּין\hdot
אֶלָּא הֲרֵי זוֹ בְּחֶזְקַת דְּרוּסַת אִישׁ עַד שֶׁתָּבִיא רְאָיָה לִדְבָרֶיהָ׃}

\halacha{רָאוּהָ מְדַבֶּרֶת עִם אֶחָד בַּשּׁוּק\hdot
אָֽמְרוּ לָהּ מַה טִּבּוֹ שֶׁל זֶה\hdot
אִישׁ פְּלוֹנִי וְכֹהֵן הוּא\hdot
רַבָּן גַּמְלִיאֵל וְרִבִּי אֱלִיעֶזֶר אוֹמְרִים נֶאֱמֶנֶת\hdot
רִבִּי יְהוֹשֻׁעַ אוֹמֵר לֹא מִפִּיהָ אָנוּ חַיִּין\hdot
אֶלָּא הֲרֵי זוֹ בְּחֶזְקַת בְּעוּלָה לְנָתִין וּלְמַמְזֵר\hdot
עַד שֶׁתָּבִיא רְאָיָה לִדְבָרֶיהָ׃}

\halacha{הָֽיְתָה מְעֻבֶּרֶת\hdot
וְאָֽמְרוּ לָהּ מַה טִּבּוֹ שֶׁל עֻבָּר זֶה\hdot
מֵאִישׁ פְּלוֹנִי וְכֹהֵן הוּא\hdot
רַבָּן גַּמְלִיאֵל וְרִבִּי אֱלִיעֶזֶר אוֹמְרִים נֶאֱמֶנֶת\hdot
רִבִּי יְהוֹשֻׁעַ אוֹמֵר לֹא מִפִּיהָ אָנוּ חַיִּין\hdot
אֶלָּא הֲרֵי זוֹ בְּחֶזְקַת מְעֻבֶּרֶת מִנָּתִין וּמַמְזֵר\hdot
עַד שֶׁתָּבִיא רְאָיָה לִדְבָרֶיהָ׃}

\halacha{אָמַר רִבִּי יוֹסֵי מַעֲשֶׂה בְּתִינוֹקֶת שֶיָּֽרְדָה לְמַלֹּאת מַיִם מִן הָעַיִן וְנֶאֱנָסָה\hdot
אָמַר רִבִּי יוֹחָנָן בֶּן נוּרִי אִם רוֹב אַנְשֵׁי הָעִיר מַשִּׂיאִים לַכְּהוּנָה הֲרֵי זוֹ תִנָּשֵׂא לַכְּהוּנָה׃}

\perek{הָאִשָּׁה}{שֶׁנִּתְאַרְמְלָה אוֹ שֶׁנִּתְגָּֽרְשָׁה\hdot
וְהִיא אוֹמֶרֶת בְּתוּלָה נְשָׂאתַנִי\hdot
וְהוּא אוֹמֵר לֹא כִי\hdot
אֶלָּא אַלְמָנָה נְשָׂאתִיךְ\hdot
אִם יֵשׁ עֵדִים שֶׁיָּצָאת בְּהִינוּמָא וְרֹאשָׁהּ פָּרוּעַ כְּתוּבָתָהּ מָאתַיִם\hdot
רִבִּי יוֹחָנָן בֶּן בְּרוֹקָא אוֹמֵר אַף חִלּוּק קְלָיוֹת רְאָיָה׃}

\halacha{וּמוֹדֶה רִבִּי יְהוֹשֻׁעַ בְּאוֹמֵר לַחֲבֵירוֹ שָׂדֶה זוֹ שֶׁל אָבִיךָ הָֽיְתָה וּלְקַחְתִּיהָ הֵימֶנּוּ\hdot
שֶׁהוּא נֶאֱמָן\hdot
שֶׁהַפֶּה שֶׁאָסַר הוּא הַפֶּה שֶׁהִתִּיר\hdot
וְאִם יֵשׁ עֵדִים שֶׁהָֽיְתָה שֶׁל אָבִיו\hdot
וְהוּא אוֹמֵר לְקַחְתִּיהָ הֵימֶנּוּ\hdot
אֵינוֹ נֶאֱמָן׃}

\halacha{הָעֵדִים שֶׁאָֽמְרוּ כְתַב יָדֵינוּ הוּא זֶה אֲבָל אֲנוּסִים הָיִינוּ\hdot
קְטַנִּים הָיִינוּ\hdot
פְּסוּלֵי עֵדוּת הָיִינוּ\hdot
הֲרֵי אֵלּוּ נֶאֱמָנִים\hdot
וְאִם יֵשׁ עֵדִים שֶׁהוּא כְתַב יָדָם\hdot
אוֹ שֶׁהָיָה כְתַב יָדָם יוֹצֵא מִמָּקוֹם אַחֵר\hdot
אֵינָן נֶאֱמָנִין׃}

\halacha{זֶה אוֹמֵר זֶה כְּתַב יָדִי\hdot
וְזֶה כְּתַב יָדוֹ שֶׁל חֲבֵרִי\hdot
וְזֶה אוֹמֵר זֶה כְּתַב יָדִי\hdot
וְזֶה כְּתַב יָדוֹ שֶׁל חֲבֵרִי\hdot
הֲרֵי אֵלּוּ נֶאֱמָנִין\hdot
זֶה אוֹמֵר זֶה כְּתַב יָדִי\hdot
וְזֶה אוֹמֵר זֶה כְּתַב יָדִי\hdot
צְרִיכִים לְצָרֵף עִמָּהֶם אַחֵר דִּבְרֵי רִבִּי\hdot
וַחֲכָמִים אוֹמְרִים אֵינָם צְרִיכִים לְצָרֵף עִמָּהֶם אַחֵר\hdot
אֶלָּא נֶאֱמָן אָדָם לוֹמַר זֶה כְּתַב יָדִי׃}

\halacha{הָאִשָּׁה שֶׁאָֽמְרָה אֵשֶׁת אִישׁ הָיִיתִי\hdot
וּגְרוּשָׁה אָנִי\hdot
נֶאֱמֶנֶת\hdot
שֶׁהַפֶּה שֶׁאָסַר הוּא הַפֶּה שֶׁהִתִּיר\hdot
וְאִם יֵשׁ עֵדִים שֶׁהָֽיְתָה אֵשֶׁת אִישׁ\hdot
וְהִיא אוֹמֶרֶת גְּרוּשָׁה אָנִי\hdot
אֵינָהּ נֶאֱמֶנֶת\hdot
אָֽמְרָה נִשְׁבֵּתִי וּטְהוֹרָה אָנִי\hdot
נֶאֱמֶנֶת\hdot
שֶׁהַפֶּה שָׁאָסַר הוּא הַפֶּה שֶׁהִתִּיר\hdot
וְאִם יֵשׁ עֵדִים שֶׁנִּשְׁבֵּית\hdot
וְהִיא אוֹמֶרֶת טְהוֹרָה אָנִי\hdot
אֵינָהּ נֶאֱמֶנֶת\hdot
וְאִם מִשֶּׁנִּשֵּׂאת בָּאוּ הָעֵדִים הֲרֵי זוֹ לֹא תֵצֵא׃}

\halacha{שְׁתֵּי נָשִׁים שֶׁנִּשְׁבּוּ\hdot
זֹאת אוֹמֶרֶת נִשְׁבֵּתִי וּטְהוֹרָה אָנִי\hdot
וְזֹאת אוֹמֶרֶת נִשְׁבֵּתִי וּטְהוֹרָה אָנִי\hdot
אֵינָן נֶאֱמָנוֹת\hdot
וּבִזְמַן שֶׁהֵן מְעִידוֹת זוֹ אֶת זוֹ\hdot
הֲרֵי אֵלּוּ נֶאֱמָנוֹת׃}

\halacha{וְכֵן שְׁנֵי אֲנָשִׁים\hdot
זֶה אוֹמֵר כֹּהֵן אָנִי\hdot
וְזֶה אוֹמֵר כֹּהֵן אָנִי\hdot
אֵינָן נֶאֱמָנִין\hdot
וּבִזְמַן שֶׁהֵן מְעִידִין זֶה אֶת זֶה\hdot
הֲרֵי אֵלּוּ נֶאֱמָנִין׃}

\halacha{רִבִּי יְהוּדָה אוֹמֵר אֵין מַעֲלִין לַכְּהוּנָה עַל פִּי עֵד אֶחָד\hdot
אָמַר רִבִּי אֶלְעָזָר אֵימָתַי בִּמְקוֹם שֶׁיֵּשׁ עוֹרְרִין\hdot
אֲבָל בִּמְקוֹם שֶׁאֵין עוֹרְרִין מַעֲלִין לַכְּהוּנָה עַל פִּי עֵד אֶחָד\hdot
רַבָּן שִׁמְעוֹן בֶּן גַּמְלִיאֵל אוֹמֵר מִשּׁוּם רִבִּי שִׁמְעוֹן בֶּן הַסְּגָן\hdot
מַעַלִין לַכְּהוּנָה עַל פִּי עֵד אֶחָד׃}

\halacha{הָאִשָּׁה שֶׁנֶּחְבְּשָׁה בִידֵי גוֹיִם עַל יְדֵי מָמוֹן\hdot
מוּתֶרֶת לְבַעְלָהּ\hdot
עַל יְדֵי נְפָשׁוֹת\hdot
אֲסוּרָה לְבַעְלָהּ\hdot
עִיר שֶׁנִּכְבְּשָׁה כַּרְכּוּם\hdot
כַָּל הַכֹּהֲנוֹת שֶׁנִּמְצְאוּ בְתוֹכָהּ\hdot
פְּסוּלוֹת\hdot
וְאִם יֵשׁ לָהֶן עֵדִים\hdot
אֲפִילוּ עֶבֶד\hdot
אֲפִילוּ שִׁפְחָה\hdot
הֲרֵי אֵלּוּ נֶאֱמָנִין\hdot
וְאֵין נֶאֱמָן אָדָם עַל יְדֵי עַצְמוֹ\hdot
אָמַר רִבִּי זְכַרְיָה בֶּן הַקָּצָב הַמָּעוֹן הַזֶּה לֹא זָזָה יָדָהּ מִתוֹךְ יָדִי מִשָּׁעָה שֶׁנִּכְנְסוּ גוֹיִם לִירוּשָׁלַיִם וְעַד שֶׁיָּֽצְאוּ\hdot
אָֽמְרוּ לוֹ אֵין אָדָם מֵעִיד עַל יְדֵי עַצְמוֹ׃}

\halacha{וְאֵלּוּ נֶאֱמָנִין לְהָעִיד בְּגָדְלָן\hdot
מַה שֶּׁרָאוּ בְקָטְנָן\hdot
נֶאֱמָן אָדָם לוֹמַר זֶה כְּתַב יָדוֹ שֶׁל אַבָּא\hdot
וְזֶה כְּתַב יָדוֹׁ שֶׁל רִבִּי\hdot
וְזֶה כְּתַב יָדוֹ שֶׁל אָחִי\hdot
זָכוּר אֲנִי בִפְלוֹנִית שֶׁיָּֽצְתָה בְּהִינוּמָא וְרֹאשָׁהּ פָּרוּעַ\hdot
וְשֶׁהָיָה אִישׁ פְּלוֹנִי יוֹצֵא מִבֵּית הַסֵּפֶר\hdot
לִטְבּוֹל לֶאֱכוֹל בַּתְּרוּמָה\hdot
וְשֶׁהָיָה חוֹלֵק עִמָּנוּ עַל הַגֹּרֶן\hdot
וְשֶׁהַמָּקוֹם הַזֶּה בֵּית הַפֶּרֶס\hdot
וְעַד כָּאן הָיִינוּ בָאִין בְּשַׁבָּת\hdot
אֲבָל אֵין אָדָם נֶאֱמָן לוֹמַר דֶּרֶךְ הָיָה לִפְלוֹנִי בַּמָּקוֹם הַזֶּה\hdot
מַעֲמָד וּמוֹשָׁב וּמִסְפֵּד\hdot
הָיָה לִפְלוֹנִי בַּמָּקוֹם הַזֶּה׃}

\chead{יום יא}

\perek{אֵלּוּ}{נַעֲרוֹת שֶׁיֵּשׁ לָהֶן קְנָס\hdot
הַבָּא עַל הַמַּמְזֶרֶת\hdot
וְעַל הַנְּתִינָה\hdot
וְעַל הַכּוּתִית\hdot
הַבָּא עַל הַגִּיוֹרֶת\hdot
וְעַל הַשְּׁבוּיָה\hdot
וְעַל הַשִּׁפְחָה\hdot
שֶׁנִּפְדּוּ וְשֶׁנִּתְגַּיְּרוּ וְשֶׁנִּשְׁתַּחְרְרוּ\hdot
פְּחוּתוֹת מִבְּנוֹת שָׁלֹשׁ שָׁנִים וְיוֹם אֶחָד\hdot
הַבָּא עַל אֲחוֹתוֹ וְעַל אֲחוֹת אָבִיו\hdot
וְעַל אֲחוֹת אִמּוֹ\hdot
וְעַל אֲחוֹת אִשְׁתּוֹ\hdot
וְעַל אֵשֶׁת אָחִיו\hdot
וְעַל אֵשֶׁת אֲחִי אָבִיו\hdot
וְעַל הַנִּדָּה\hdot
יֵשׁ לָהֶן קְנָס\hdot
אַף עַל פִּי שֶׁהֵן בִּהִכָּרֵת אֵין בָּהֶן מִיתַת בֵּית דִּין׃}

\halacha{וְאֵלּוּ שֶׁאֵין לָהֶן קְנָס\hdot
הַבָּא עַל הַגִּיוֹרֶת\hdot
וְעַל הַשְּׁבוּיָה\hdot
וְעַל הַשִּׁפְחָה\hdot
שֶׁנִּפְדּוּ וְשֶׁנִּתְגַּיְּרוּ וְשֶׁנִּשְׁתַּחְרְרוּ\hdot
יְתֵרירוֹת עַל בְּנוֹת שָׁלֹשׁ שָׁנִים וְיוֹם אֶחָד\hdot
רִבִּי יְהוּדָה אוֹמֵר שְׁבוּיָה שֶׁנִּפְדֵּית הֲרֵי הִיא  בִקְדוּשָׁתָהּ\hdot
אַף עַל פִּי שֶׁגְּדוֹלָה\hdot
הַבָּא עַל בִּתּוֹ\hdot
וְעַל בַּת בִּתּוֹ\hdot
וְעַל בַּת בְּנוֹ\hdot
וְעַל בַּת אִשְׁתּוֹ\hdot
וְעַל בַּת בְּנָהּ\hdot
וְעַל בַּת בִּתָּהּ\hdot
אֵין לָהֶן קְנָס\hdot
מִפְּנֵי שֶׁמִּתְחַיֵּיב בְּנַפְשׁוֹ\hdot
שֶׁמִּיתָתוֹ בִּידֵי בֵית דִּין\hdot
וְכָל הַמִּתְחַיֵּיב בְּנַפְשׁוֹ אֵינוֹ מְשַׁלֵּם מָמוֹן\hdot
שֶׁנֶּאֱמַר אִם לֹא יִהְיֶה אָסוֹן עָנוֹש יֵעָנֵשׁ׃}

\halacha{נַעֲרָה שֶׁנִּתְאָֽרְסָה\hdot
וְנִתְגָּֽרְשָׁה\hdot
רִבִּי יוֹסֵי הַגְּלִילִי אוֹמֵר אֵין לָהּ קְנָס\hdot
רִבְּי עֲקִיבָא אוֹמֵר יֵשׁ לָהּ קְנָס וּקְנָסָהּ לְעַצְמָהּ׃}

\halacha{הַמְּפַתֶּה נוֹתֵן שְׁלֹשָׁה דְבָרִים\hdot
הָאוֹנֵס אַרְבָּעָה\hdot
הַמְּפַתֶּה נוֹתֵן בֹּשֶׁת וּפְגָם וּקְנָס\hdot
מוֹסִיף עָלָיו הָאוֹנֵס\hdot
שֶׁנּוֹתֵן אֶת הַצַּעַר\hdot
מַה בֵּין הָאוֹנֵס לַמְּפַתֶּה\hdot
הָאוֹנֵס נותֵן אֶת הַצַּעַר\hdot
הַמְּפַתֶּה אֵינוֹ נוֹתֵן אֶת הַצַּעַר\hdot
הָאוֹנֵס נוֹתֵן מִיָּד\hdot
וְהַמְּפַתֶּה לִכְשֶׁיּוֹצִיא\hdot
הָאוֹנֵס שׁוֹתֶה בַעֲצִיצוֹ\hdot
וְהַמְּפַתֶּה אִם רָצָה לְהוֹצִיא מוֹצִיא׃}

\halacha{כֵּיצַד שׁוֹתֶה בַעֲצִיצוֹ\hdot
אֲפִילוּ הִיא חִגֶּרֶת\hdot
אֲפִילוּ הִיא סוֹמָא\hdot
וַאֲפִילוּ הִיא מֻכַּת שְׁחִין\hdot
נִמְצָא בָהּ דְּבַר עֶרְוָה\hdot
אוֹ שֶׁאֵינָהּ רְאוּיָה לָבֹא בְיִשְׂרָאֵל\hdot
אֵינוֹ רַשָּׁאי לְקַיְּמָהּ\hdot
שֶׁנֶּאֱמַר וְלוֹ תִהְיֶה לְאִשָּׁה אִשָּׁה הָֽרְאוּיָה לוֹ׃}

\halacha{יְתוֹמָה שֶׁנִּתְאָֽרְסָה וְנִתְגָּֽרְשָׁה\hdot
רִבִּי אֶלְעָזָר אוֹמֵר הַמְּפַתֶּה פָּטוּר\hdot
וְהָאוֹנֵס חַיָּיב׃}

\halacha{אֵיזֶהוּ בֹּשֶׁת\hdot
הַכֹּל לְפִי הַַמְּבַיֵּשׁ וְהַמִּתְבַּיֵּשׁ\hdot
פְּגָם רוֹאִין אוֹתָהּ כְּאִלּוּ הִיא שִׁפְחָה נִמְכֶּרֶת\hdot
כַּמָּה הָֽיְתָה יָפָה\hdot
וְכַמָּה הִיא יָפָה\hdot
קְנָס שָׁוֶה בְּכָל אָדָם\hdot
וְכָל שֶׁיֵּשׁ לוֹ קִצְבָּה מִן הַתּוֹרָה שָׁוֶה בְּכָל אָדָם׃}

\halacha{כָּל מָקוֹם שֶׁיֵּשׁ מֶכֶר אֵין קְנָס\hdot
וְכָל מָקוֹם שֶׁיֵּשׁ קְנָס אֵין מֶכֶר\hdot
קְטַנָּה יֵשׁ לָהּ מֶכֶר וְאֵין לָהּ קְנָס\hdot
נַעֲרָה יֵשׁ לָהּ קְנָס וְאֵין לָהּ מֶכֶר\hdot
הַבּוֹגֶרֶת אֵין לָהּ לֹא מֶכֶר וְלֹא קְנָס׃}

\halacha{הָאוֹמֵר פִּתִּיתִי אֶת בִּתּוֹ שֶׁל פְּלוֹנִי\hdot
מְשַׁלֵּם בֹּשֶׁת וּפְגָם עַל פִּי עַצְמוֹ\hdot
וְאֵינוֹ מְשַׁלֵּם קְנָס\hdot
הָאוֹמֵר גָּנַבְתִּי (וְטָבַחְתִּי וּמָכַרְתִּי)\hdot
מְשַׁלֵּם אֶת הַקֶּרֶן עַל פִּי עַצְמוֹ\hdot
וְאֵינוֹ מְשַׁלֵּם תַּשְׁלומֵי כֶפֶל\hdot
וְתַשְׁלוּמֵי אַרְבָּעָה וַחֲמִשָּׁה\hdot
הֵמִית שׁוֹרִי אֶת פְּלוֹנִי\hdot
אוֹ שׁוֹרוֹ שֶׁל פְּלוֹנִי\hdot
הֲרֵי זֶה מְשַׁלֵּם עַל פִּי עַצְמוֹ\hdot
הֵמִית שׁוֹרִי עַבְדּוֹ שֶׁל פְּלוֹנִי\hdot
אֵינוֹ מְשַׁלֵּם עַל פִּי עַצְמוֹ\hdot
זֶה הַכְּלָל\hdot
כָּל הַמְּשַׁלֵּם יֶתֶר עַל מַה שֶׁהִזִּיק אֵינוֹ מְשַׁלֵּם עַל פִּי עַצְמוֹ׃}

\perek{נַעֲרָה}{שֶׁנִּתְפַּתְּתָה\hdot
בָּשְׁתָּהּ וּפְגָמָהּ וּקְנָסָהּ שֶׁל אָבִיהָ\hdot
וְהַצַּעַר בַּתְּפוּסָה\hdot
עָֽמְדָה בְדִין עַד שֶׁלֹּא מֵת הָאָב\hdot
הֲרֵי הֵן שֶׁל אָב\hdot
מֵת הָאָב\hdot
הֲרֵי הֵן שֶׁל אָחִים\hdot
לֹא הִסְפִּיקָה לַעֲמוֹד בְּדִין עַד שֶׁמֵּת הָאָב\hdot
הֲרֵי הֵן שֶׁל עַצְמָהּ\hdot
רִבִּי שִׁמְעוֹן אוֹמֵר אִם לֹא הִסְפִּיקָה לְגַבּוֹת עַד שֶׁמֵּת הָאָב\hdot
הֲרֵי הֵן שֶׁל עַצְמָהּ\hdot
עָֽמְדָה בְדִין עַד שֶׁלֹא בָֽגְרָה\hdot
הֲרֵי הֵן שֶׁל אָב\hdot
מֵת הָאָב\hdot
הֲרֵי הֵן שֶׁל אָחִים\hdot
לֹא הִסְפִּיקָה לַעֲמוֹד בְּדִין עַד שֶׁבָּֽגְרָה\hdot
הֲרֵי הֵן שֶׁל עַצְמָהּ\hdot
רִבִּי שִׁמְעוֹן אוֹמֵר אִם לֹא הִסְפִּיקָה לְגַבּוֹת עַד שֶׁמֵּת הָאָב\hdot
הֲרֵי הֵן שֶׁל עַצְמָהּ\hdot
מַעֲשֶׂה יָדֶיהָ וּמְצִיאָתָהּ אַף עַל פִּי שֶׁלֹּא גָֽבְתָה\hdot
מֵת הָאָב\hdot
הֲרֵי הֵן שֶׁל אָחִים׃}

\halacha{הַמְּאָרֵס אֶת בִּתּוֹ וְגֵרְשָׁהּ\hdot
אֵרְסָהּ וְנִתְאַרְמְלָה כְּתוּבָתָהּ שֶׁלּוֹ\hdot
הִשִּׂיאָהּ וְגֵרְשָׁהּ\hdot
הִשִּׂיאָהּ וְנִתְאַרְמְלָה כְּתוּבָתָּהּ שֶׁלָּהּ\hdot
רִבִּי יְהוּדָה אוֹמֵר הָרִאשׁוֹנָה שֶׁל אָב\hdot
אָֽמְרוּ לוֹ (מִשֶּׁהִשִּׂיאָהּ) אֵין לְאָבִיהָ רְשׁוּת בָּהּ׃}

\halacha{הַגִּיּוֹרֶת שֶׁנִּתְגַּיְּרָה בִּתָּהּ עִמָּהּ וְזִנְּתָה\hdot
הֲרֵי זוֹ בְחָנֵק\hdot
אֵין לָהּ לֹא פֶתַח בֵּית הָאָב\hdot
וְלֹא מֵאָה סֶלַע\hdot
הָֽיְתָה הוֹרָתָהּ שֶׁלֹּא בִקְדוּשָׁה וְלֵידָתָהּ בִּקְדוּשָׁה\hdot
הֲרֵי זוֹ בִסְקִילָה\hdot
אֵין לָהּ לֹא פֶתַח בֵּית הָאָב וְלֹא מֵאָה סֶלַע\hdot
הָֽיְתָה הוֹרָתָהּ וְלֵידָתָהּ בִּקְדוּשָׁה\hdot
הֲרֵי הִיא כְבַת יִשְׂרָאֵל לְכָל דָּבָר\hdot
יֵשׁ לָהּ אָב וְאֵין לָהּ פֶּתַח בֵּית הָאָב\hdot
יֵשׁ לָהּ פֶּתַח בֵּית הָאָב וְאֵין לָהּ אָב\hdot
הֲרֵי זוֹ בִסְקִילָה\hdot
לֹא נֶאֱמַר פֶּתַח בֵּית אָבִיהָ אֶלָּא לְמִצְוָה׃}

\halacha{הָאָב זַכָּאי בְּבִתּוֹ בִּקִדּוּשֶׁיהָ\hdot
בְּכֶסֶף בִּשְׁטַר וּבְבִיאָה\hdot
וְזַכָּאי בִּמְצִיאָתָהּ וּבְמַעֲשֵׂה יָדֶיהָ וּבַהֲפָרַת נְדָרֶיהָ\hdot
וּמְקַבֵּל אֶת גִּטָּהּ\hdot
וְאֵינוֹ אוֹכֵל פֵּרוֹת בְּחַיֵֶּיהָ\hdot
נִשֵּׂאת\hdot
יָתֵר עָלָיו הַבַּעַל שֶׁאוֹכֵל פֵּירוֹת בְּחַיֶּיהָ וְחַיָּיב בִּמְזוֹנוֹתֶיהָ\hdot
בְּפִרְקוּנָהּ וּבִקְבוּרָתָהּ\hdot
רִבִּי יְהוּדָה אוֹמֵר אֲפִילוּ עָנִי שֶׁבְּיִשְׂרָאֵל\hdot
לֹא יִפְחוֹת מִשְּׁנֵי חֲלִילִים וּמְקוֹנֶנֶת׃}

\halacha{לְעוֹלָם הִיא בִּרְשׁוּת הָאָב עַד שֶׁתִּכָּנֵס לִרְשׁוּת הַבַּעַל לַנִּשּׂוּאִין\hdot
מָסַר הָאָב לְשִׁלּוּחֵי הַבַּעַל\hdot
הֲרֵי הִיא בִּרְשׁוּת הַבַּעַל\hdot
הַָלַךְ הָאָב עִם שִׁלּוּחֵי הַבַּעַל\hdot
אוֹ שֶׁהָֽלְכוּ שִׁלּוּחֵי הָאָב עִם שִׁלּוּחֵי הַבַּעַל\hdot
הֲרֵי הִיא בִּרְשׁוּת הָאָב\hdot
מָֽסְרוּ שִׁלּוּחֵי הָאָב לְשִׁלּוּחֵי הַבַּעַל\hdot
הֲרֵי הִיא בִּרְשׁוּת הַבַּעַל׃}

\halacha{הָאָב אֵינוֹ חַיָּיב בִּמְזוֹנוֹת בִּתּוֹ\hdot
זֶה מִדְרָשׁ דָּרַשׁ רִבִּי אֶלְעָזָר בֶּן עֲזַרְיָה לִפְנֵי חֲכָמִים בַּכֶּרֶם בְּיַבְנֶה\hdot
הַבָּנִים יִירְשׁוּ וְהַבָּנוֹת יִזּוֹנוּ\hdot
מַה הַבָּנִים אֵינָם יוֹרְשִׁים\hdot
אֶלָּא לְאַחַר מִיתַת אֲבִיהֶן\hdot
אַף הַבָּנוֹת אֵינָן נִזּוֹנוֹת אֶלָּא לְאַחַר מִיתַת אֲבִיהֶן׃}

\halacha{לֹא כָתַב לָהּ כְּתוּבָה\hdot
בְּתוּלָה גוֹבָה מָאתַיִם\hdot
וְאַלְמָנָה מָנֶה\hdot
מִפְּנֵי שֶׁהוּא תְּנָאי בֵּית דִּין\hdot
כָּתַב לָהּ שָׂדֶה שָׁוֶה מָנֶה תַּחַת מָאתַיִם\hdot
וְלֹא כָתַב לָהּ כָּל נְכָסִים דְּאִית לִי אַחֲרָאִין לִכְתוּבָתִיךְ\hdot
חַיָּיב שֶׁהוּא תְּנָאי בֵּית דִּין׃}

\halacha{לֹא כָתַב לָהּ\hdot
אִם תִּשְׁתַבָּאי אֲפָֽרְקִינָךְ וְאוֹתְבִינָךְ לִי לְאִנְתּוּ\hdot
וְבַכֹּהֶנֶת אֲהַדְרִינָךְ לִמְדִינָתִיךְ חַיָּיב שֶׁהוּא תְּנָאי בֵּית דִּין׃}

\halacha{נִשְׁבֵּית חַַיָּיב לִפְדוֹתָהּ\hdot
וְאִם אָמַר הֲרֵי גִטָּהּ וּכְתוּבָתָהּ תִּפְדֶּה אֶת עַצְמָהּ\hdot
אֵינוֹ רַשָּׁאי\hdot
לָֽקְתָה חַיָּיב לְרַפֹּאתָהּ\hdot
אָמַר הֲרֵי גִטָּהּ וּכְתוּבָתָהּ תְּרַפֵּא אֶת עַצְמָהּ רַשָּׁאי׃}

\halacha{לֹא כָתַב לָהּ בְּנִין דִּכְרִין דִּיהֲווֹ לִיכִי מִנָּאי\hdot
אִינוּן יִרְתוּן כֶּסֶף כְּתוּבָתִיךְ\hdot
יָתֵר עַל חוֹלַקְהוֹן דְּעִם אֲחוּהוֹן\hdot
חַיָּיב שֶׁהוּא תְּנָאי בֵּית דִּין׃}

\halacha{בְּנָן נוּקְבָן דִּיהַוְיָן לִיכִי מִנָּאי\hdot
יְהַוְיָן יַתְבָן בְּבֵיתִי\hdot
וּמִתְזָנָן מִנְּכָסַי עַד דְּתִנְסְבָן לְגֻבְרִין\hdot
חַיָּיב שֶׁהוּא תְּנָאי בֵּית דִּין׃}

\halacha{אַתְּ תְּהֵא יַתְבָא בְּבֵיתִי וּמִתְזָנָא מִנְּכָסַי\hdot
כָּל יְמֵי מִיגַד אַלְמְנוּתִיךְ בְּבֵיתִי\hdot
חַיָּיב שֶׁהוּא תְּנָאי בֵּית דִּין\hdot
כַּך הָיוּ אַנְשֵׁי יְרוּשָׁלַיִם כּוֹתְבִין\hdot
וְאַנְשֵׁי גָלִיל הָיוּ כוֹתְבִין כְּאַנְשֵׁי יְרוּשָׁלַיִם\hdot
אַנְשֵׁי יְהוּדָה הָיוּ כוֹתְבִין\hdot
עַד שֶׁיִּרְצוּ הַיּוֹרְשִׁין לִתֵּן לִיךְ כְּתוּבָתִיךְ\hdot
לְפִיכַךְ אִם רָצוּ יוֹרְשִׁין נוֹתְנִין לָהּ כְּתוּבָתָהּ\hdot
וּפוֹטְרִין אוֹתָהּ׃}

\perek{אַף}{עַל פִּי שֶׁאָֽמְרוּ בְּתוּלָה גוֹבָה מָאתַיִם וְאַלְמָנָה מָנֶה\hdot
אִם רָצָה לְהוֹסִיף אֲפִילוּ מֵאָה מָנֶה יוֹסִיף\hdot
נִתְאַרְמְלָה אוֹ נִתְגָּֽרְשָׁה\hdot
בֵּין מִן הָאֵרוּסִין בֵּין מִן הַנִּשּׂוּאִין\hdot
גּוֹבָה אֶת הַכֹּל\hdot
רִבִּי אֶלְעָזָר בֶּן עֲזַרְיָה אוֹמֵר מִן הַנִּשּׂוּאִין גּוֹבָה אֶת הַכֹּל\hdot
מִן הָאֵרוּסִין בְּתוּלָה גוֹבָה מָאתִים\hdot
וְאַלְמָנָה מָנֶה\hdot
שֶׁלֹּא כָתַב לָהּ אֶלָּא עַל מְנַת לְכוֹנְסָהּ\hdot
רִבִּי יְהוּדָה אוֹמֵר אִם רָצָה כוֹתֵב לַבְּתוּלָה שְׁטַר שֶׁל מָאתַיִם\hdot
וְהִיא כוֹתֶבֶת הִתְקַבַּלְתִּי מִמָּךְ מָנֶה\hdot
וּלְאַלְמָנָה מָנֶה\hdot
וְהִיא כוֹתֶבֶת הִתְקַבַּלְתִּי מִמָּךְ חֲמִשִּׁים זוּז\hdot
רִבִּי מֵאִיר אוֹמֵר כָּל הַפּוֹחֵת לַבְּתוּלָה מִמָּאתַיִם\hdot
וְלָאַלְמָנָה מִמָּנֶה\hdot
הֲרֵי זוֹ בְּעִילַת זְנוּת׃}

\halacha{נוֹתְנִין לַבְּתוּלָה שְׁנֵים עָשָׂר חֹדֶשׁ מִשֶּׁתְּבָעָהּ הַבַּעַל לְפַרְנֵס אֶת עַצְמָהּ\hdot
וּכְשֵׁם שֶׁנּוֹתְנִין לָאִשָּׁה\hdot
כַּךְ נוֹתְנִין לָאִישׁ לְפַרְנֵס אֶת עַצְמוֹ\hdot
וְלָאַלְמָנָה שְׁלֹשִׁים יוֹם\hdot
הִגִּיעַ זְמָן וְלֹא נִשְּׂאוּ אוֹכְלוֹת מִשֶּׁלּוֹ\hdot
וְאוֹכְלוֹת בַּתְּרוּמָה\hdot
רִבִּי טַרְפוֹן אוֹמֵר נוֹתְנִין לָהּ הַכֹּל תְּרוּמָה\hdot
רִבִּי עֲקִיבָא אוֹמֵר מֶחְצָה חוֹלִין וּמֶחְצָה תְרוּמָה׃}

\halacha{הַיָּבָם אֵינוֹ מַאֲכִיל בַּתְּרוּמָה\hdot
עָֽשְׂתָה שִׁשָּׁה חֳדָשִׁים בִּפְנֵי הַבַּעַל\hdot
וְשִׁשָּׁה חֳדָשִׁים בִּפְנֵי הַיָּבָם\hdot
וַאֲפִילוּ כֻלָּן בִּפְנֵי הַבַּעַל\hdot
חָסֵר יוֹם אֶחָד בִּפְנֵי הַיָּבָם\hdot
אוֹ כֻלָּן בִּפְנֵי הַיָּבָם\hdot
חָסֵר יוֹם אֶחָד בִּפְנֵי הַבַּעַל אֵינָהּ אוֹכֶלֶת בַּתְּרוּמָה\hdot
זוֹ מִשְׁנָה רִאשׁוֹנָה\hdot
בֵּית דִּין שֶׁל אַחֲרֵיהֶן אָֽמְרוּ אֵין הָאִשָּׁה אוֹכֶלֶת בַּתְּרוּמָה עַד שֶׁתִּכָּנֵס לַחֻפָּה׃}

\halacha{הַמַּקְדִּישׁ מַעֲשֵׂה יְדֵי אִשְׁתּוֹ\hdot
הֲרֵי זוֹ עוֹשָׂה וְאוֹכֶלֶת\hdot
הַמּוֹתָר רִבִּי מֵאִיר אוֹמֵר הֶקְדֶּשׁ\hdot
רִבִּי יוֹחָנָן הַסַּנְדְּלָר אוֹמֵר חוֹלִין׃}

\halacha{אֵלּוּ מְלָאכוֹת שֶׁהָאִשָּׁה עוֹשָׂה לְבַעְלָהּ\hdot
טוֹחֶנֶת וְאוֹפָה וּמְכַבֶּסֶת וּמְבַשֶּׁלֶת וּמְנִיקָה אֶת בְּנָהּ\hdot
וּמַצַּעַת לוֹ הַמִּטָּה\hdot
ועוֹשָׂה בַצֶּמֶר\hdot
הִכְנִיסָה לוֹ שִׁפְחָה אַחַת\hdot
לֹא טוֹחֶנֶת\hdot
וְלֹא אוֹפָה\hdot
וְלֹא מְכַבֶּסֶת\hdot
שְׁתַּיִם אֵינָהּ מְבַשֶּׁלֶת\hdot
וְאֵינָהּ מְנִיקָה אֶת בְּנָהּ\hdot
שָׁלֹשׁ אֵינָהּ מַצַּעַת (לוֹ) הַמִּטָּה\hdot
וְאֵינָהּ עוֹשָׂה בַצֶּמֶר\hdot
אַרְבַּע יוֹשֶׁבֶת בְּקַתֶּדְרָא\hdot
רִבִּי אֱלִיעֶזֶר אוֹמֵר אֲפִילוּ הִכְנִיסָה לוֹ מֵאָה שְׁפָחוֹת כּוֹפָהּ לַעֲשׂוֹת בַּצֶּמֶר\hdot
שֶׁהַבַּטָּלָה מְבִיאָה לִידֵי זִמָּה\hdot
רַבָּן שִׁמְעוֹן בֶּן גַּמְלִיאֵל אוֹמֵר אַף הַמַּדִּיר אֶת אִשְׁתּוֹ מִלַּעֲשׂוֹת מְלָאכָה\hdot
יוֹצִיא וְיִתֵּן כְּתוּבָה\hdot
שֶׁהַבַּטָּלָה מְבִיאָה לִידִי שִׁיעֲמוּם׃}

\halacha{הַמַדִּיר אֶת אִשְׁתּוֹ מִתַּשְׁמִישׁ הַמִּטָּה\hdot
בֵּית שַׁמָּאי אוֹמְרִים שְׁתֵּי שַׁבָּתוֹת\hdot
בֵּית הִלֵּל אוֹמְרִים שַׁבָּת אֶחָת\hdot
הַתַּלְמִידִים יוֹצְאִים לְתַלְמוּד תּוֹרָה שֶׁלֹּא בִרְשׁוּת שְׁלֹשִים יוֹם\hdot
הַפּוֹעֲלִים שַׁבָּת אֶחָת\hdot
הָעוֹנָה הָאֲמוּרָה בַתּוֹרָה\hdot
הַטַּיָּלִין בְּכָל יוֹם\hdot
הַפּוֹעֲלִים שְׁתַּיִם בְּשַׁבָּת\hdot
הַחַמָּרִים אַחַת בְּשַׁבָּת\hdot
הַגַּמָּלִים אֶחָד לִשְׁלֹשִׁים יוֹם\hdot
הַסַּפָּנִים אַחַת לְשִׁשָּׁה חֳדָשִׁים דִּבְרֵי רִבִּי אֱלִיעֶזֶר׃}

\halacha{הַמּוֹרֶדֶת עַל בַּעְלָהּ\hdot
פּוֹחֲתִין לָהּ מִכְּתוּבָתָהּ שִׁבְעָה דִינָרִין בְּשַׁבָּת\hdot
רִבִּי יְהוּדָה אוֹמֵר שִׁבְעָה טַרְפָעִיקִים\hdot
עָד מָתַי הוּא פוֹחֵת\hdot
עַד כְּנֶגֶד כְּתוּבָתָהּ\hdot
רִבִּי יוֹסֵי אוֹמֵר לְעוֹלָם הוּא פוֹחֵת וְהוֹלֵךְ שֶׁמֶּא תִפּוֹל לָהּ יְרוּשָׁה מִמָּקוֹם אַחֵר\hdot
וְיַחֲזוֹר וְיִגְבֶּה הֵימֶנָּה\hdot
וְכֵן הַמּוֹרֵד עַל אִשְׁתּוֹ מוֹסִיפִין לָהּ עַל כְּתוּבָתָהּ שְׁלֹשָה דִינָרִים בְּשַׁבָּת\hdot
רִבִּי יְהוּדָה אוֹמֵר שְׁלֹשָׁה טַרְפָעִיקִין׃}

\halacha{הַמְּשָׁרֶה אֶת אִשְׁתּוֹ עַל יְדֵי שָׁלִישׁ לֹא יִפְחוֹת לָהּ מִשְּׁנֵי קַבִּין חִטִּין\hdot
אוֹ מֵאַרְבַּע קַבִּין שְׂעוֹרִים\hdot
אָמַר רִבִּי יוֹסֵי לֹא פָסַק לָהּ שְׂעוֹרִים אֶלָּא רִבִּי יִשְׁמָעֵאל שֶׁהָיָה סָמוּךְ לֶאֱדוֹם\hdot
וְנוֹתֵן לָהּ חֲצִי קָב קָטְנִית\hdot
וַחֲצִי לוֹג שֶׁמֶן\hdot
וְקָב גְּרוֹגְרוֹת אוֹ מָנֶה דְבֵילָה\hdot
וְאִם אֵין לוֹ\hdot
פּוֹסֵק לְעֻמָּתָן פֵּירוֹת מִמָּקוֹם אַחֵר\hdot
וְנוֹתֵן לָהּ מִטָּה וּמַפֵּץ וּמַחֲצֶלֶת\hdot
וְנוֹתֵן לָהּ כִּפָּה לְרֹאשָׁהּ\hdot
וַחֲגוֹר לְמָתְנֶיהָ\hdot
וּמַנְעָלִים מִמּוֹעֵד לְמוֹעֵד\hdot
וְכֵלִים שֶׁל חֲמִשִּׁים זוּז מִשָּׁנָה לְשָׁנָה\hdot
וְאֵינוֹ נוֹתֵן לָהּ לֹא חֲדָשִׁים בִּימוֹת הַחַמָּה\hdot
וְלֹא שְׁחָקִים בִּימוֹת הַגְּשָׁמִים\hdot
אֶלָּא נוֹתֵן לָהּ כֵּלִים שֶׁל חֲמִשִּׁים זוּז בִּימוֹת הַגְּשָׁמִים\hdot
וְהִיא מִתְכַּסָּה בְּבַלָּאוֹתֵיהֶן בִּימוֹת הַחַמָּה\hdot
וְהַשְּׁחָקִים שֶׁלָּהּ׃}

\halacha{נוֹתֵן לָהּ מְעָה כֶסֶף לְצָרְכֶיהָ\hdot
וְאוֹכֶלֶת עִמּוֹ מִלֵּילֵי שַׁבָּת לְלֵילֵי שַׁבָּת\hdot
וְאִם אֵין נוֹתֵן לָהּ מְעָה כֶסֶף לְצָרְכֶיהָ\hdot
מַעֲשֵׂה יָדֶיהָ שֶׁלָּהּ\hdot
וּמָה הִיא עוֹשָׂה לוֹ\hdot
מִשְׁקַל חֲמֵֵשׁ סְלָעִים שְׁתִי בִּיהוּדָה\hdot
שֶׁהֵן עֶשְׂרִים סְלָעִים בַּגָּלִיל\hdot
אוֹ מִשְׁקַל עֶשֶׂר סְלָעִים עֶרֶב בִּיהוּדָה\hdot
שֶׁהֵן עֶשְׂרִים סְלָעִים בַּגָּלִיל\hdot
וְאִם הָֽיְתָה מְנִיקָה פּוֹחֲתִין לָהּ מִמַּעֲשֵׂה יָדֶיהָ\hdot
וּמוֹסִיפִין לָהּ עַל מְזוֹנוֹתֶיהָ\hdot
בַּמֶּה דְּבָרִים אֲמוּרִים בֶּעָנִי שֶׁבְּיִשְׂרָאֵל\hdot
אֲבָל בַּמֻּכְבָּד הַכֹּל לְפִי כְבוֹדוֹ׃}

\perek{מְצִיאַת}{הָאִשָּׁה וּמַעֲשֵׂה יָדֶיהָ לְבַעְלָהּ\hdot
וִירוּשָׁתָהּ הוּא אוֹכֵל פֵּירוֹת בְּחַיֶּיהָ\hdot
בָּשְׁתָּהּ וּפְגָמָהּ\hdot
שֶׁלָּהּ\hdot
רִבִּי יְהוּדָה בֶן בְּתֵירָא אוֹמֵר\hdot
בִּזְמַן שֶׁבַּסֵּתֶר לָהּ שְׁנֵי חֲלָקִים וְלוֹ חֵלֶק אֶחָד\hdot
וּבִזְמַן שֶׁבַּגָּלוּי\hdot
לוֹ שְׁנֵי חֲלָקִים וְלָהּ אֶחָד\hdot
שֶׁלּוֹ יִנָּתֵן מִיָּד\hdot
וְשֶׁלָּהּ יִלָּקַח בָּהֶן קַרְקַע וְהוּא אוֹכֵל פֵּירוֹת׃}

\halacha{הַפּוֹסֵק מָעוֹת לֲחֲתָנוֹ וּמֵת חֲתָנוֹ\hdot
אָֽמְרוּ חֲכָמִים יָכוֹל הוּא שֶׁיֹּאמַר\hdot
לְאָחִיךָ הָיִיתִי רוֹצֶה לִתֵּן\hdot
וְלָךְ אִי אַפְשִׁי לִתֵּן׃}

\halacha{פָּֽסְקָה לְהַכְנִיס לוֹ אֶלֶף דִּינָר\hdot
הוּא פוֹסֶק כְּנֶגְדָּן חֲמִשָּׁה עָשָׂר מָנֶה\hdot
וּכְנֶגֶד הַשּׁוֹם\hdot
הוּא פוֹסֵק פָּחוּת חֹמֶשׁ\hdot
שׁוֹם בְּמָנֶה וְשָׁוֶה מָנֶה אֵין לוֹ אֶלָּא מָנֶה\hdot
שׁוֹם בְּמָנֶה הִיא נוֹתֶנֶת שְׁלֹשִׁים וְאֶחָד סֶלַע וְדִינָר\hdot
וּבְאַרְבַּע מֵאוֹת הִיא נוֹתֶנֶת חֲמֵשׁ מֵאוֹת\hdot
מַה שֶּׁחָתָן פּוֹסֵק הוּא פוֹסֵק פָּחוּת חֹמֶשׁ׃}

\halacha{פָּֽסְקָה לְהַכְנִיס לוֹ כְּסָפִים סֶלַע (כֶּסֶף) נַעֲשָׂה שִׁשָּׁה דִינָרִים\hdot
הֶחָתָן מְקַבֵּל עָלָיו עֲשָׂרָה דִינָרִין לַקֻּפָּה לְכָל מָנֶה וּמָנֶה\hdot
רַבָּן שִׁמְעוֹן בֶּן גַּמְלִיאֵל אוֹמֵר הַכֹּל כְּמִנְהַג הַמְּדִינָה׃}

\halacha{הַמַּשִּׂיא אֶת בִּתּוֹ סְתָם\hdot
לֹא יִפְחוֹת לָהּ מֵחֲמִשִּׁים זוּז\hdot
פָּסַק לְהַכְנִיסָהּ עֲרוּמָה לֹא יֹאמַר הַבַּעַל כְּשֶׁתָּבֹא לְבֵיתִי אֲכַסֶּנָּה בִכְסוּתִי\hdot
אֶלָּא מְכַסָּהּ וְעוֹדָהּ בְּבֵית אָבִיהָ\hdot
וְכֵן הַמַּשִּׂיאִין אֶת הַיְּתוֹמָה לֹא יִפְחֲתוּ לָהּ מֵחֲמִשִּׁים זוּז\hdot
אִם יֵשׁ בְּכִיס מְפַרְנְסִין אוֹתָהּ לְפִי כְבוֹדָהּ׃}

\halacha{יְתוֹמָה שֶׁהִשִּׂיאוּהָ אִמָּהּ אוֹ אַחֶיהָ מִדַּעְתָּהּ\hdot
וְכָֽתְבוּ לָהּ בְּמֵאָה אוֹ בַחֲמִשִּׁים זוּז\hdot
יְכוֹלָה הִיא מִשֶּׁתַּגְדִּיל\hdot
לְהוֹצִיא מִיָּדָן מַה שֶּׁרָאוּי לְהִנָּתֵן לָהּ\hdot
רִבִּי יְהוּדָה אוֹמֵר אִם הִשִּׂיא אֶת הַבַּת הָרִאשׁוֹנָה יְנָּתֵן לַשְּׁנִיָּה כְּדֶרֶךְ שֶׁנָּתַן לָרִאשׁוֹנָה\hdot
וַחֲכָמִים אוֹמְרִים פְּעָמִים שֶׁאָדָם עָנִי וְהֶעֱשִׁיר\hdot
אוֹ עָשִׁיר וְהֶעֱנִי\hdot
אֶלָּא שָׁמִין אֶת הַנְּכָסִים וְנוֹתְנִין לָהּ׃}

\halacha{הַמַּשְׁלִישׁ מָעוֹת לְבִתּוֹ\hdot
וְהִיא אוֹמֶרֶת נֶאֱמָן בַּעֲלִי עָלַי\hdot
יַעֲשֶׂה הַשָּׁלִישׁ מַה שֶּׁהֻשְׁלַשׁ בְּיָדוֹ\hdot
דִּבְרֵי רִבִּי מֵאִיר\hdot
רִבִּי יוֹסֵי אוֹמֵר וְכִי אֵינָהּ אֶלָּא שָׂדֶה\hdot
הִיא רוֹצָה לְמוֹכְרָהּ\hdot
הֲרֵי הִיא מְכוּרָה מֵעַכְשָׁיו\hdot
בַּמֶּה דְּבָרִים אֲמוּרִים בַּגְּדוֹלָה\hdot
אֲבָל בַּקְּטַנָּה אֵין מַעֲשֵׁה קְטַנָּה כְּלוּם׃}

\perek{הַמַּדִּיר}{אֶת אִשְׁתּוֹ מִלִּיהֲנוֹת לוֹ\hdot
עַד שְׁלֹשִׁים יוֹם\hdot
יַעֲמִיד פַּרְנָס\hdot
יָתֶר מִכֵּן יוֹצִיא וְיִתֵּן כְּתוּבָה\hdot
רִבִּי יְהוּדָה אוֹמֵר\hdot
בְּיִשְׁרָאֵל חֹדֶשׁ אֶחָד יְקַיֵּם\hdot
וּשְׁנַיִם יוֹצִיא וְיִתֵּן כְּתוּבָה\hdot
וּבַכֹּהֶנֶת שְׁנַיִם יְקַיֵּם\hdot
וּשְׁלֹשָׁה יוֹצִיא וְיִתֵּן כְּתוּבָה׃}

\tractatehead{מסכת סוטה}{sota}

\perek{הַמְּקַנֵּא}{לְאִשְׁתּוֹ\hdot
רִבִּי אֱלִיעֶזֶר אוֹמֵר מְקַנֵּא לָהּ עַל פִּי שְׁנַיִם\hdot
וּמַשְׁקָהּ  עַל פִּי עֵד אֶחָד\hdot
אוֹ עַל פִּי עַצְמוֹ\hdot
רִבִּי יְהוֹשֻׁעַ אוֹמֵר מְקַנֵּא לָהּ עַל פִּי שְׁנַיִם וּמַשְׁקָהּ עַל פִּי שְׁנַיִם׃}

\halacha{כֵּיצַד מְקַנֵּא לָהּ\hdot
אָמַר לָהּ בִּפְנֵי שְׁנַיִם\hdot
אַל תְּדַבְּרִי עִם אִישׁ פְּלוֹנִי\hdot
וְדִבְּרָה עִמּוֹ\hdot
עֲדַיִן הִיא מֻתֶּרֶת לְבֵיתָהּ\hdot
וּמֻתֶּרֶת לֶאֱכוֹל בַּתְּרוּמָה\hdot
נִכְנְסָה עִמּוֹ לְבֵית הַסֵּתֶר\hdot
וְשָׁהֲתָה עִמּוֹ כְּדֵי טֻמְאָה\hdot
אֲסוּרָה לְבֵיתָהּ וַאֲסוּרָה לֶאֱכוֹל בַּתְּרומָה\hdot
וְאִם מֵת חוֹלֶצֶת וְלֹא מִתְיַבֶּמֶת׃}

\halacha{וְאֵלּוּ אֲסוּרוֹת מִלֶּאֱכוֹל בַּתְּרוּמָה\hdot
הָאוֹמֶרֶת טְמֵאָה אֲנִי לָךְ\hdot
וְשֶׁבָּאוּ עֵדִים שֶׁהִיא טְמֵאָה\hdot
וְהָאוֹמֶרֶת אֲנִי שׁוֹתָה\hdot
וְשֶׁבַּעְלָהּ אֵינוֹ רוֹצֶה לְהַשְׁקוֹתָהּ\hdot
וְשֶׁבַּעְלָהּ בָּא עֲלֶיהָ בַּדֶּרֶךְ\hdot
כֵּיצַד עוֹשֶׂה לָהּ\hdot
מוֹלִיכָהּ לְבֵית דִּין שֶׁבְּאוֹתוֹ מָקוֹם\hdot
וּמוֹסְרִין לוֹ שְׁנֵי תַלְמִידֵי חֲכָמִים שֶׁמֶא יָבֹא עָלֶיהָ בַּדֶּרֶךְ\hdot
רִבִּי יְהוּדָה אוֹמֵר בַּעְלָהּ נֶאֱמָן עָלֶיהָ׃}

\halacha{הָיוּ מַעֲלִין אוֹתָהּ לְבֵית דִּין הַגָּדוֹל שֶׁבִּירוּשָׁלַיִם\hdot
וּמְאַיְּמִין עָלֶיהָ כְּדֶרֶךְ שֶׁמְּאַיְּמִין עַל עֵדֵי נְפָשׁוֹת\hdot
וְאוֹמְרִים לָהּ בִּתִּי הַרְבֶּה יַיִן עוֹשֶׂה\hdot
הַרְבֶּה שְׂחוֹק עוֹשֶׂה\hdot
הַרְבֶּה יַלְדּוּת עוֹשָׂה\hdot
הַרְבֶּה שְׁכֵנִים הָרָעִים עוֹשִׂים\hdot
עֲשִׂי לִשְׁמוֹ הַגָּדוֹל שֶׁנִּכְתַּב בִּקְדֻשָּׁה שֶׁלֹּא יִמָּחֶה עַל הַמַּיִם\hdot
וְאוֹמְרִים לְפָנֶיהָ דְּבָרִים שֶׁאֵינָהּ כְּדַי לְשׁוֹמְעָן\hdot
הִיא וְכָל מִשְׁפַּחַת בֵּית אָבִיהָ׃}

\halacha{אִם אָֽמְרָה טְמֵאָה אָנִי\hdot
שׁוֹבֶרֶת כְּתוּבָתָהּ וְיוֹצֵאת\hdot
וְאִם אָֽמְרָה טְהוֹרָה אָנִי\hdot
מַעֲלִין אוֹתָהּ לְשַׁעַר הַמִּזְרַח\hdot
שֶׁעַל פֶּתַח שַׁעַר נִקָּנוֹר\hdot
שֶׁשָּׁם מַשְׁקִין אֶת הַסּוֹטוֹת\hdot
וּמְטַהֲרִין אֶת הַיּוֹלְדוֹת\hdot
וּמְטַהֲרִין אֶת הַמְּצוֹרָעִים\hdot
וְכֹהֵן אוֹחֵז  בִּבְגָדֶיהָ\hdot
אִם נִקְרָעוּ נִקְרָעוּ\hdot
וְאִם נִפְרָמוּ נִפְרָמו\hdot
עַד שֶׁהוּא מְגַלֶּה אֶת לִבָּהּ\hdot
וְסוֹתֵר אֶת שְׂעָרָהּ\hdot
רִבִּי יְהוּדָה אוֹמֵר אִם הָיָה לִבָּהּ נָאֶה\hdot
לֹא הָיָה מְגַלֵּהוּ\hdot
ואִם הָיָה שְׁעָרָהּ נָאֶה\hdot
לֹא הָיָה סוֹתְרוֹ׃}

\halacha{הָֽיְתָה מְכֻּסָּה בִלְבָנִים\hdot
מְכַסָּהּ בִּשְׁחְוֹרִים\hdot
הָיָה עָלֶיהָ כְּלֵי זָהָב וְקַטְלִיאוֹת\hdot
נְזָמִים וְטַבָּעוֹת\hdot
מַעֲבִירָם מִמֶּנָּה כְּדֵי לְנַוְּלָהּ\hdot
וְאַחַר כַּךְ מֵבִיא חֶבֶל מִצְרִי וְקוֹשֵׁר לְמַעְלָה מִדַּדֶּיהָ\hdot
וְכָל הָרוֹצֶה לִרְאוֹת\hdot
בָּא לִרְאוֹת חוּץ מֵעֲבָדֶיהָ וְשִׁפְחוֹתֶיהָ\hdot
מִפְּנֵי שֶׁלִּבָּהּ גַּס בָּהֶן\hdot
וְכָל הַנָּשִׁים מֻתָּרוֹת לִרְאוֹתָהּ\hdot
שֶׁנֶּאֱמַר וְנִוָּֽסְרוּ כָל הַנָּשִׁים וְלֹא תַעֲשֶׂינָה כּזִמַתְכֶנָה׃}

\halacha{בַּמִּדָּה שֶׁאָדָם מוֹדֵד בָּהּ מוֹדְדִין לוֹ\hdot
הִיא קִשְּׁטָה אֶת עַצְמָהּ לַעֲבֵירָה\hdot
הַמָּקוֹם נִוְּלָהּ\hdot
הִיא גִלְּתָה אֶת עַצְמָהּ לַעֲבֵירָה\hdot
הַמָּקוֹם גִּלָּה עָלֶיהָ\hdot
הַיָּרֵךְ הִתְחִילָה בַעֲבֵירָה תְּחִלָּה\hdot
וְאַחַר כַּךְ הַבָּטֶן\hdot
וּשְׁאַר כָּל הַגּוּף לֹא פָלַט׃}

\halacha{שִׁמְשׁוֹן הָלַךְ אַחַר עֵינָיו\hdot
לְפִיכַךְ נָֽקְרוּ פְלִשְׁתִּים אֶת עֵינָיו\hdot
שֶׁנֶּאֱמַר וַיֹּאחֲזוּהוּ פְלִשְׁתִּים וַיְּנַקְּרוּ אֶת עֵינָיו\hdot
אַבְשָׁלוֹם נִתְגָּאָה בְשַׂעֲרוֹ\hdot
לְפִיכַךְ נִתְלָה בְשַׂעֲרוֹ\hdot
וּלְפִי שֶׁבָּא עַל עֶשֶׂר פִּלַגְשֵׁי אָבִיו\hdot
לְפִיכַךְ נָֽתְנוּ בוֹ עֶשֶׂר לֻנְבִיּוֹת\hdot
שֶׁנֶּאֱמַר וַיָּסֹבּוּ עֲשָׂרָה נְעָרִים נוֹשְׂאֵי כְלֵי יוֹאָב\hdot
וּלְפִי שֶׁגָּנַב  שְׁלֹשָׁה לְבָבוֹת\hdot
לֵב אָבִיו\hdot
וְלֵב בֵּית  דִּין\hdot
ולֵב יִשְׂרָאֵל\hdot
שֶׁנֶּאֱמַר וַיְּגַנֵּב אַבְשָׁלוֹם אֶת לֵב אַנְשֵׁי יִשְׂרָאֵל\hdot
לְפִיכַךְ נִתְקְעוּ בוֹ שְׁלֹשָׁה שְׁבָטִים\hdot
שֶׁנֶּאֱמַר וַיִּקַח שְׁלֹשָׁה שְׁבָטִים בְּכַפּוֹ וַיִּתְקָעֵם בְּלֵב אַבְשָׁלוֹם׃}

\halacha{וְכֵן לְעִנְיַן הַטּוֹבָה\hdot
מִרְיָם הִמְתִּינָה לְמֹשֶׁה שָׁעָה אַחַת\hdot
שֶׁנֶּאֱמַר וַתֵּתַצַּב אֲחוֹתוֹ מֵרָחוֹק\hdot
לְפִיכַךְ נִתְעַכְּבוּ לָהּ יִשְׁרָאֵל שִׁבְעָה יָמִים בַּמִּדְבָּר שֶׁנֶּאֱמַר וְהָעָם לֹא נָסַע עַד הֵאָסֵף מִרְיָם\hdot
יוֹסֵף זָכָה לִקְבּוֹר אֶת אָבִיו\hdot
וְאֵין בְּאֶחָיו גָּדוֹל מִמֶּנּוּ\hdot
שֶׁנֶּאֱמַר וַיַּעַל יוֹסֵף לִקְבּוֹר אֶת אָבִיו וְגוֹמֵר\hdot
וַיַּעַל עִמּוֹ גַּם רֶכֶב וְגַם פָּרָשִׁים\hdot
מִי לָנוּ גָּדוֹל מִיּוֹסֵף\hdot
שֶׁלֹּא נִתְעַסַּק בּוֹ אֶלָּא מֹשֶׁה\hdot
מֹשֶׁה זָכָה  בְּעַצְמוֹת יוֹסֵף\hdot
וְאֵין בְּיִשְׂרָאֵל  גָּדוֹל מִמֶּנּוּ\hdot
שֶׁנֶּאֱמַר וַיִּקַּח מֹשֶׁה אֶת עַצְמוֹת יוֹסֵף עִמּוֹ\hdot
מִי לָנוּ גָּדוֹל מִמֹּשֶׁה שֶׁלֹּא נִתְעָסַק בּוֹ אֶלָּא הַמָּקוֹם\hdot
שֶׁנֶּאֱמַר וַיִּקְבּוֹר אוֹתוֹ בַxגַּיְא\hdot
לֹא עַל מֹשֶׁה  בִלְבַד אָֽמְרוּ\hdot
אֶלָּא עַל כָּל הַצַּדִּיקִים\hdot
שֶׁנֶּאֱמַר וְהָלַךְ  לְפָנֶיךָ צִדְקֶךָ כְּבוֹד יְיָ יַאַסְפֶךָ׃}

\tractatehead{מסכת נדרים}{nedarim}

\perek{כָּל}{כִּנּוּיֵי נְדָרִים כִּנְדָרִים\hdot
וַחֲרָמִים כַּחֲרָמִים\hdot
וּשְׁבוּעוֹת כִּשְׁבוּעוֹת\hdot
וּנְזִירוֹת כִּנְזִירוֹת\hdot
הָאוֹמֵר לַחֲבֵירוֹ מוּדַרְנִי מִמָּךְ\hdot
מֻפְרַשְׁנִי מִמָּךְ\hdot
מְרוּחַקְנִי מִמָּךְ\hdot
שֶׁאֵינִי אוֹכֵל לָךְ\hdot
שֶׁאֵינִי טוֹעֵם לָךְ אָסוּר\hdot
מְנוּדֶה אֲנִי לָךְ\hdot
רִבִּי עֲקִיבָא הָיָה חוֹכֵךְ בַּזֶּה לְהַחְמִיר\hdot
כְּנִדְרֵי רְשָׁעִים נָדַר בַּנָּזִיר וּבְקָרְבָּן וּבִשְׁבוּעָה\hdot
כְּנִדְרֵי כְשֵׁרִים\hdot
לֹא אָמַר כְּלוּם\hdot
כְּנִדְבוֹתָם נָדַר בַּנָּזִיר וּבַקָּרְבָּן׃}

\halacha{הָאוֹמֵר לַחֲבֵירוֹ\hdot
קוּנָם\hdot
קוּנָח\hdot
קוּנָס\hdot
הֲרֵי אֵלּוּ כִּנּוּיִן לַקָּרְבָּן\hdot
חֵרֶק\hdot
חֵרֶךְ\hdot
חֵרֶף\hdot
הֲרֵי אֵלּוּ כִּנּוּיִם לַחֵרֶם\hdot
נָזִיק\hdot
נָזִיחַ\hdot
פָּזִיחַ\hdot
הֲרֵי אֵלּוּ כִּנּוּיִן לִנְזִירוּת\hdot
שְׁבוּתָה\hdot
שְׁקוּקָה\hdot
נָדַר בְּמוֹתָה\hdot
הֲרֵי אֵלּוּ כִּנּוּיִן לַשְּׁבוּעָה׃}

\halacha{הָאוֹמֵר לֹא חוֹלִין לֹא אוֹכַל לָךְ\hdot
לֹא כָשֵׁר וְלֹא דָכֵי\hdot
טָהוֹר\hdot
וְטָמֵא\hdot
נוֹתָר\hdot
וּפִגּוּל\hdot
אָסוּר\hdot
כְּאִימְרָא\hdot
כְּדִירִים\hdot
כְּעֵצִים\hdot
כְּאִשִּׁים\hdot
כְּמִזְבֵּחַ\hdot
כְּהֵיכָל\hdot
כִּירוּשָׁלַיִם\hdot
נָדַר בְּאַחַד מִכָּל מְשַׁמְּשֵׁי הַמִּזְבֵּחַ\hdot
אַף עַל פִּי שֶׁלֹּא הִזְכִּיר קָרְבָּן\hdot
הֲרֵי זֶה נָדַר קָרְבָּן\hdot
רִבִּי יְהוּדָה אוֹמֵר הָאוֹמֵר יְרוּשָׁלַיִם לֹא אָמַר כְּלוּם׃}

\halacha{הָאוֹמֵר קָרְבָּן\hdot
עוֹלָה\hdot
מִנְחָה\hdot
חַטָּאת\hdot
תּוֹדָה\hdot
שְׁלָמִים\hdot
שֶׁאֵינִי אוֹכֵל לָךְ\hdot
אָסוּר\hdot
רִבִּי יְהוּדָה מַתִּיר\hdot
הַקָּרְבָּן\hdot
כְּקָרְבָּן\hdot
קָרְבָּן שֶאוֹכַל לָךְ\hdot
אָסוּר\hdot
לְקָרְבָּן לֹא אוֹכַל לָךְ\hdot
רִבִּי מֵאִיר אוֹסֵר\hdot
הָאוֹמֵר לַחֲבֵירוֹ קוּנָם פִּי מְדַבֵּר עִמָּךְ\hdot
יָדִי עוֹשָׂה עִמָּךְ\hdot
רַגְלִי מְהַלֶּכֶת עִמָּךְ\hdot
אָסוּר׃}

\perek{וְאֵלּוּ}{מֻתָּרִין חוֹלִין שֶׁאוֹכַל לָךְ\hdot
כִּבְשַׂר חֲזִיר\hdot
כַּעֲבוֹדַת כּוֹכָבִים\hdot
כְּעוֹרוֹת לְבוּבִין\hdot
כִּנְבֵלוֹת\hdot
(כִּטְרֵפוֹת)\hdot
כִּשְׁקָצִים\hdot
(כִּרְמָשִׁים)\hdot
כְּחַלַּת אַהֲרֹן וְכִתְרוּמָתוֹ מֻתָּר\hdot
הָאוֹמֵר לְאִשְׁתּוּ הֲרֵי אַתְּ עָלַי כְּאִמָּא\hdot
פּוֹתְחִין לוֹ פֶּתַח מִמָּקוֹם אַחֵר\hdot
שֶׁלֹּא יָקֵל רֹאשׁוֹ לְכַךְ\hdot
קוּנָם שֶׁאֵינִי יָשֵׁן\hdot
שֶׁאֵינִי מְדַבֵּר\hdot
שֶׁאֵינִי מְהַלֵּךְ\hdot
הָאוֹמֵר לְאִשָּׁה קוּנָם שֶׁאֵינִי מְשַׁמְּשָׁךְ\hdot
הֲרֵי זֶה בְּלֹא יָחֵל דְּבָרוֹ\hdot
שְׁבוּעָה שֶׁאֵינִי יָשֵׁן\hdot
שֶׁאֵינִי מְדַבֵּר\hdot
שֶׁאֵינִי מְהַלֵּךְ אָסוּר׃}

\halacha{קָרְבָּן לֹא אוֹכַל לָךְ\hdot
הַקָּרְבָּן שֶׁאוֹכַל לָךְ\hdot
לֹא קָרְבָּן לֹא אוֹכַל לָךְ\hdot
מֻתָּר\hdot
שְׁבוּעָה לֹא אוֹכַל לָךְ\hdot
שְׁבוּעָה שֶׁאוֹכַל לָךְ\hdot
לֹא שְׁבוּעָה לֹא אוֹכַל לָךְ\hdot
אָסוּר\hdot
זֶה חֹמֶר בַּשְּׁבוּעוֹת מִבַּנְּדָרִים\hdot
וְחֹמֶר בַּנְּדָרִים מִבַּשְּׁבוּעוֹת\hdot
כֵּיצַד\hdot
אָמַר קוּנָם סֻכָּה שֶׁאֵינִי עוֹשֶׂה\hdot
לוּלָב שֶׁאֵינִי נוֹטֵל\hdot
תְּפִילִין שֶׁאֵינִי מַנִּיחַ\hdot
בַּנְּדָרִים אָסוּר\hdot
בַּשְּׁבוּעוֹת מֻתָּר\hdot
שֶׁאֵין נִשְׁבָּעִין לַעֲבוֹר עַל הַמִּצְוֺת׃}

\halacha{יֵשׁ נֶדֶר בְּתוֹךְ נֶדֶר\hdot
וְאֵין שְׁבוּעָה בְּתוֹךְ שְׁבוּעָה\hdot
כֵּיצַד\hdot
אָמַר הֲרֵינִי נָזִיר אִם אוֹכַל\hdot
הֲרֵינִי נָזִיר אִם אוֹכַל וְאָכַל\hdot
חַיָּיב עַל כָּל אַחַת וְאֶחָת\hdot
שְבוּעָה שֶׁלֹּא אוֹכַל\hdot
שְׁבוּעָה שֶׁלֹּא אוֹכַל וְאָכַל\hdot
אֵינוֹ חַיָּיב אֶלָּא אֶחָת׃}

\halacha{סְתָם נְדָרִים לְהַחְמִיר וּפֵירוּשָׁם לְהָקֵל\hdot
כֵּיצַד\hdot
אָמַר הֲרֵי עָלַי כִּבְשַׂר מָלִיחַ\hdot
כְּיַיִן נֶסֶךְ\hdot
אִם שֶׁל שָׁמַיִם נָדַר\hdot
אָסוּר\hdot
וְאִם שֶׁל עֲבוֹדַת אֱלִילִים מֻתָּר\hdot
וְאִם סְתָם אָסוּר\hdot
הֲרֵי עָלַי כְּחֵרֶם אִם כְּחֵרֶם שֶׁל שָׁמַיִם אָסוּר\hdot
וְאִם כְּחֵרֶם שֶׁל כֹּהֲנִים מֻתָּר\hdot
וְאִם סְתָם אָסוּר\hdot
הֲרֵי עָלַי כְּמַעֲשֵׂר אִם כְּמַעְשַׂר בְּהֵמָה נָדַר אָסוּר\hdot
וְאִם שֶׁל דָּגָן מֻתָּר\hdot
וְאִם סְתָם אָסוּר\hdot
הֲרֵי עָלַי כִּתְרוּמָה\hdot
אְם כִּתְרוּמַת הַלִּשְׁכָּה נָדַר אָסוּר\hdot
וְאִם שֶׁל גּוֹרֶן מֻתָּר\hdot
וְאִם סְתָם אָסוּר\hdot
דִּבְרֵי רִבִּי מֵאִיר\hdot
רִבִּי יְהוּדָה אוֹמֵר סְתָם תְּרוּמָה בִּיהוּדָה אֲסוּרָה\hdot
וּבַגָּלִיל מוּתֶרֶת\hdot
שֶׁאֵין אַנְשֵׁי גָלִיל מַכִּירִין אֶת תְּרוּמַת הַלִּשְׁכָּה\hdot
סְתָם חֲרָמִים בִּיהוּדָה מֻתָּרִין וּבַגָּלִיל אֲסוּרִין\hdot
שֶׁאֵין אַנְשֵׁי גָלִיל מַכִּירִין אֶת חֶרְמֵי הַכֹּהֲנִים׃}

\halacha{נָדַר בְּחֵרֶם\hdot
וְאָמַר לֹא נָדַרְתִּי אֶלָּא בְּחֶרְמוֹ שֶׁל יָם\hdot
בְּקָרְבָּן וְאָמַר לֹא נָדַרְתִּי אֶלָּא בְּקָרְבְּנוֹת שֶׁל מְלָכִים\hdot
הֲרֵי עַצְמִי קָרְבָּן\hdot
וְאָמַר לֹא נָדַרְתִּי אֶלָּא בְּעֶצֶם שֶׁהִנַּחְתִּי לִי לִהְיוֹת נוֹדֵר בּוֹ\hdot
קוּנָם אִשְׁתִּי נֶהֱנֵית לִי\hdot
וְאָמַר לֹא נָדַרְתִּי אֶלָּא מֵאִשְׁתִּי הָרִאשׁוֹנָה שֶׁגֵּרַשְׁתִּי\hdot
עַל כֻּלָּן אֵין נִשְׁאָלִים לָהֶם\hdot
וְאִם נִשְׁאֲלוּ עוֹנְשִׁין אוֹתָן\hdot
וּמַחְמִירִין עֲלֵיהֶם\hdot
דִּבְרֵי רִבִּי מֵאִיר\hdot
וַחֲכָמִים אוֹמְרִים פּוֹתְחִין לָהֶן פֶּתַח מִמָּקוֹם אַחֵר וּמְלַמְּדִין אוֹתָם כְּדֵי שֶׁלֹּא יְנַהֲגוּ קַלּוּת רֹאשׁ בַּנְּדָרִים׃}

\perek{אַרְבָּעָה}{נְדָרִים הִתִּירוּ חֲכָמִים\hdot
נִדְרֵי זֵרוּזִין\hdot
וְנִדְרֵי הֲבָאי\hdot
וְנִדְרֵי שְׁגָגוֹת\hdot
וְנִדְרֵי אוֹנָסִים\hdot
נִדְרֵי זֵירוּזִין\hdot
כֵּיצַד\hdot
הָיָה מוֹכֵר חֵפֶץ\hdot
וְאָמַר קוּנָם שֶׁאֵינִי פוֹחֵת לָךְ מִן הַסֶּלַע\hdot
וְהַלָּה אוֹמֵר קוּנָם שֶׁאֵינִי מוֹסִיף לָךְ עַל הַשֶׁקֶל\hdot
שְׁנֵיהֶם רוֹצִים בִּשְׁלֹשָׁה דִינָרִין\hdot
רִבִּי אֱלִיעֶזֶר בֶּן יַעֲקֹב אוֹמֵר אַף הָרוֹצֶה לְהַדִּיר אֶת חֲבֵירוֹ שֶׁיֹּאכַל אֶצְלוֹ\hdot
אוֹמֵר (כָּל) נֶדֶר שֶׁאֲנִי עָתִיד לִדּוֹר הוּא בָטֵל\hdot
וּבִלְבָד שֶׁיְּהֵא זָכוּר בְּשַׁעַת הַנֶּדֶר׃}

\halacha{נִדְרֵי הֲבָאי\hdot
אָמַר קוּנָם אִם לֹא רָאִיתִי בַּדֶּרֶךְ הַזֶּה כְּיוֹצְאֵי מִצְרַיִם\hdot
אִם לֹא רָאִיתִי נָחָשׁ כְּקוֹרַת בֵּית הַבָּד\hdot
נִדְרֵי שְׁגָגוֹת אִם אָכַלְתִּי וְאִם שָׁתִיתִי\hdot
וְנִזְכַּר שֶׁאָכַל וְשָׁתָה\hdot
שֶׁאֵינִי אוֹכֵל וְשֶׁאֵינִי שׁוֹתֶה\hdot
וְשָׁכַח וְאָכַל וְשָׁתָה\hdot
אָמַר קוּנָם אִשְׁתִּי נֶהֱנֵית לִי שֶׁגָֽנְבָה אֶת כִּיסִי וְשֶׁהִכְּתָה אֶת בְּנִי\hdot
וְנוֹדַע שֶׁלֹּא הִכַּתּוּ\hdot
וְנוֹדַע שֶׁלֹּא גְנָבַתּוּ\hdot
רָאָה אוֹתָן אוֹכְלִים תְּאֵנִים\hdot
וְאָמַר הֲרֵי הֵן עֲלֵיכֶם קָרְבָּן\hdot
וְנִמְצְאוּ אָבִיו וְאֶחָיו וְהָיוּ עִמָּהֶן אֲחֵרִים\hdot
בֵּית שַׁמָּאי אוֹמְרִים הֵן מֻתָּרִין\hdot
וּמַה שֶּׁעִמָּהֶם אֲסוּרִים\hdot
וּבֵית הִלֵּל אוֹמְרִים אֵלּוּ וְאֵלּוּ מֻתָּרִין׃}

\halacha{נִדְרֵי אוֹנָסִים\hdot
הִדִּירוֹ חֲבֵירוֹ שֶׁיֹּאכַל אֶצְלוֹ\hdot
וְחָלָה הוּא אוֹ שֶׁחָלָה בְנוֹ\hdot
אוֹ שֶׁעִכְּבוֹ נָהָר\hdot
הֲרֵי אֵלּוּ נִדְרֵי אוֹנָסִין׃}

\halacha{נוֹדְרִין לַהוֹרְגִים\hdot
וְלַחֲרָמִין\hdot
וְלַמּוֹכְסִין\hdot
שֶׁהִיא תְרוּמָה אַף עַל פִּי שֶׁאֵינָהּ תְּרוּמָה\hdot
שֶׁהֵן שֶׁל בֵּית הַמֶּלֶךְ\hdot
אַף עַל פִי שֶׁאֵינָן שֶׁל בֵּית הַמֶּלֶךְ\hdot
בֵּית שַׁמָּאי אוֹמְרִים בְּכָל נוֹדְרִין חוּץ מִבַּשְּׁבוּעָה\hdot
ובֵית הִלֵּל אוֹמְרִים אַף בַּשְּׁבוּעָה\hdot
בֵּית שַׁמָּאי אוֹמְרִים לֹא יִפְתַח לוֹ בַנֶּדֶר\hdot
ובֵית הִלֵּל אוֹמְרִים אַף יִפְתַח לוֹ\hdot
בֵּית שַׁמָּאי אוֹמְרִים בְּמַה שֶּׁהוּא מַדִּירוֹ\hdot
וּבֵית הִלֵּל אוֹמְרִים אַף בְּמַה שֶּׁאֵינוֹ מַדִּירוֹ\hdot
כֵּיצַד\hdot
אָֽמְרוּ לוֹ אֱמוֹר קוּנָם אִשְׁתִּי נֶהֱנֵית לִי\hdot
וְאָמַר קוּנָם אִשְׁתִּי וּבָנַי נֶהֱנִין לִי\hdot
בֵּית שַׁמָּאי אוֹמְרִים אִשְׁתּוֹ מֻתֶּרֶת וּבָנָיו אֲסוּרִין\hdot
וּבֵית הִלֵּל אוֹמְרִים אֵלּוּ וְאֵלּוּ מֻתָּרִין׃}

\halacha{הֲרֵי נְטִיעוֹת הָאֵלּוּ קָרְבָּן\hdot
אִם אֵינָן נִקְצָצוֹת\hdot
טָלֵית זוֹ קָרְבָּן אִם אֵינָהּ נִשְׂרֶפֶת\hdot
יֵשׁ לָהֶן פִּדְיוֹן\hdot
הֲרֵי נְטִיעוֹת הָאֵלּוּ קָרְבָּן עַד שֶׁיִּקָּצֵצוּ\hdot
טָלֵית זוֹ קָרְבָּן עַד שֶׁתִּשָּׂרֵף אֵין לָהֶן פִּדְיוֹן׃}

\halacha{הַנּודֵר מִיּוֹרְדֵי הַיָּם מֻתָּר בְּיוֹשְׁבֵי הַיַּבָּשָה\hdot
מִיּוֹשְׁבֵי הַיַּבָּשָׁה\hdot
אָסוּר בְּיוֹרְדֵי הַיָּם\hdot
שֶׁיּוֹרְדֵי הַיָּם בִּכְלַל יוֹשְׁבֵי הַיַּבָּשָׁה\hdot
לֹא כְאֵלּוּ שֶׁהוֹלְכִין מֵעַכּוֹ לְיָפוֹ אֶלָּא בְּמִי שֶׁדַּרְכּוֹ לְפָרֵשׁ׃}

\halacha{הַנּוֹדֵר מֵרוֹאֵי הַחַמָּה\hdot
אָסוּר אַף בַּסּוֹמִין\hdot
שֶׁלֹּא נִתְכַּוָּן זֶה אֶלָּא לְמִי שֶׁהַחַמָּה רוֹאָה אוֹתוֹ׃}

\halacha{הַנּוֹדֵר מִשְּׁחוֹרֵי הָרֹאשׁ אָסוּר בַּקְּרָחִין וּבַעֲלֵי שֵׂיבוֹת\hdot
וּמֻתָּר בַּנָּשִׁים וּבַקְּטַנִּים\hdot
שֶׁאֵין נִקְרָאִין שְׁחוֹרֵי הָרֹאשׁ אֶלָּא אֲנָשִׁים׃}

\halacha{הַנּוֹדֵר מִן הַיִּלּוֹדִים מֻתָּר בַּנּוֹלָדִים\hdot
(מִן הַנּוֹלָדִים אָסוּר בַּיִּלּוֹדִים)\hdot
רִבִּי מֵאִיר מַתִּיר אַף בַּיִּלּוֹדִים\hdot
וַחֲכָמִים אוֹמְרִים לֹא נִתְכַּוָּן זֶה אֶלָּא בְמִי שֶׁדַּרְכּוֹ לְהוֹלִיד׃}

\halacha{הַנּוֹדֵר מִשּׁוֹבְתֵי שַׁבָּת\hdot
אָסוּר בְּיִשְׁרָאֵל\hdot
(וְאָסוּר) בַּכּוּתִים\hdot
מֵאוֹכְלֵי שׁוּם אָסוּר בְּיִשְׂרָאֵל וְאָסוּר בַּכּוּתִים\hdot
מֵעוֹלֵי יְרוּשָׁלַיִם\hdot
אָסוּר בְּיִשְׂרָאֵל\hdot
וּמֻתָּר בַּכּוּתִים׃}

\halacha{קוּנָם שֶׁאֵינִי נֶהֱנֶה לִבְנֵי נֹחַ\hdot
מֻתָּר בְּיִשְׂרָאֵל וְאָסוּר בָּאוּמוֹת\hdot
שֶׁאֵינִי נֶהֱנֶה לְזֶרַע אַבְרָהָם\hdot
אָסוּר בְּיִשְׂרָאֵל וּמֻתָּר בָּאוּמוֹת\hdot
שֶׁאֵינִי נֶהֱנֶה לְיִשְׂרָאֵל\hdot
לוֹקֵחַ בְּיוֹתֵר וּמוֹכֵר בְּפָחוּת\hdot
שֶׁיִשְׂרָאֵל נֶהֱנִין לִי\hdot
לוֹקֵחַ בְּפָחוּת וּמוֹכֵר בְּיוֹתֵר\hdot
(אִם שׁוֹמְעִין לוֹ)\hdot
שֶאֵינִי נֶהֱנֶה לָהֶן וְהֵן לִי\hdot
יֶהֱנֶה לָאוּמוֹת\hdot
קוּנָם שֶׁאֵינִי נֶהֱנֶה לָעֲרֵלִים מֻתָּר בְּעַרְלֵי יִשְׂרָאֵל\hdot
וְאָסוּר בְּמוּלֵי גוֹיִם\hdot
קוּנָם שֶׁאֵינִי נֶהֱנֶה לַמּוּלִים\hdot
אָסוּר בְּעַרְלֵי יִשְׂרָאֵל\hdot
וּמֻתָּר בְּמוּלֵי אוּמוֹת הָעוֹלָם\hdot
שֶׁאֵין הָעָרְלָה קְרוּיָה אֶלָּא לְשֵׁם גּוֹיִם\hdot
שֶׁנֶּאֱמַר כִּי כָּל הַגּוֹיִם עֲרֵלִים וְכָל בֵּית יִשְׁרָאֵל עַרְלֵי לֵב\hdot
וְאוֹמֵר וְהָיָה הַפְּלִשְׁתִּי הֶעָרֵל הַזֶּה\hdot
וְאוֹמֵר פֶּן תִּשְׂמַחְנָה בְּנוֹת פְּלִשְׁתִּים פֶּן תַּעֲלוֹזְנָה בְּנוֹת הָעֲרֵלִים\hdot
(רִבִּי אֶלְעָזָר בֶּן עֲזַרְיָה אוֹמֵר מְאוּסָה הָעָרְלָה שֶׁנִּתְגַּנּו בָּה הָרְשָׁעִים\hdot
שֶׁנֶּאֱמַר כִּי כָּל הַגּוֹיִם עֲרֵלִים וְכָל בֵּית יִשְׁרָאֵל עַרְלֵי לֵב)\hdot
רִבִּי יִשְׁמָעֵאל אוֹמֵר גְּדוֹלָה מִילָה שֶׁנִּכְרְתוּ עָלֶיהָ שְׁלֹשׁ עֶשְׂרֵה בְּרִיתוֹת\hdot
רִבִּי יוֹסֵי אוֹמֵר גְּדוֹלָה מִילָה שֶׁדּוֹחָה אֶת הַשַּׁבָּת הַחֲמוּרָה\hdot
רִבִּי יְהוֹשֻׁעַ בֶּן קָרְחָה אוֹמֵר גְּדוֹלָה מִילָה שֶׁלֹּא נִתְלָה לוֹ לְמֹשֶׁה הַצַּדִּיק עָלֶיהָ מְלֹא שָׁעָה\hdot
רִבִּי נְחֶמְיָא אוֹמֵר גְּדוֹלָה מִילָה שֶׁדּוֹחָה אֶת הַנְּגָעִים\hdot
רִבִּי אוֹמֵר גְּדוֹלָה מִילָה שֶׁכָּל הַמִצְוֺת שֶׁעָשָׂה אַבְרָהָם אָבִינוּ לֹא נִקְרָא שָׁלֵם עַד שֶׁמָּל\hdot
שֶׁנֶּאֱמַר הִתְהַלֵּךְ לְפָנַי וֶהְיֵה תָמִים\hdot
דָּבָר אַחֵר גְּדוֹלָה מִילָה שֶׁאִלְמָלֵא הִיא לא בָרָא הַקָּדוֹשׁ בָּרוּךְ הוּא אֶת עוֹלָמוֹ\hdot
שֶׁנֶּאֱמַר כֹּה אָמַר יְיָ אִם לֹא בְרִיתִי יוֹמָם וָלַיְלָה חֻקּוֹת שָׁמַיִם וָאָרֶץ לֹא שַׂמְתִּי׃}

\chead{יום יג}

\perek{אֵין}{בֵּין הַמֻּדָּר הֲנָאָה מֵחֲבֵירוֹ\hdot
לְמֻדָּר הֵימֶנּוּ מַאֲכָל\hdot
אֶלָּא דְרִיסַת הָרֶגֶל וְכֵלִים שֶׁאֵין עוֹשִׂין בָּהֶן אֹכֶל נֶפֶשׁ\hdot
הַמֻּדָּר מַאֲכָל מֵחֲבֵירוֹ\hdot
לֹא יַשְׁאִילֶנּוּ  נָפָה וּכְבָרָה וְרֵיחַיִם וְתַנּוּר\hdot
אֲבָל מַשְׁאִיל לוֹ חַלּוּק וְטַבַּעַת\hdot
וְטָלֵית וּנְזָמִים\hdot
וְכָל דָּבָר שֶׁאֵין עוֹשִׂים בּוֹ אֹכֶל נֶפֶשׁ\hdot
מָקוֹם שֶׁמַּשְׂכִּירִין כַּיּוֹצֵא בָהֶן אָסוּר׃}

\halacha{הַמֻּדָּר הֲנָאָה מֵחֲבֵירוֹ\hdot
שׁוֹקֵל אֶת שִׁקְלוֹ\hdot
ופוֹרֵע אֶת חוֹבוֹ
וּמַחֲזִיר לוֹ אֶת אֲבֵידָתוֹ\hdot
מָקוֹם שֶׁנּוֹטְלִים עָלֶיהָ שָׂכָר תִּפּוֹל הֲנָאָה לַהֶקְדֵּשׁ׃}

\halacha{וְתוֹרֵם אֶת תְּרוּמָתוֹ וּמַעְשְׂרוֹתָיו לְדַעְתּוֹ\hdot
וּמַקְרִיב עָלָיו קִנֵּי זָבִין\hdot
קִנֵּי זָבוֹת\hdot
קִנֵּי יוֹלְדוֹת\hdot
חַטָּאוֹת וַאֲשָׁמוֹת\hdot
וּמְלַמְּדוֹ מִדְרָשׁ הֲלָכוֹת וַאֲגָדוֹת\hdot
(אֲבָל) לֹא יְלַמְּדֶנּוּ מִקְרָא\hdot
אֲבָל מְלַמֵּד הוּא אֶת בָּנָיו (וְאֶת בְּנוֹתָיו) מִקְרָא\hdot
וְזָן אֶת אִשְׁתּוֹ וְאֶת בָּנָיו\hdot
אַף עַל פִּי שֶׁהוּא חַיָּב בִּמְזוֹנוֹתָיו\hdot
וְלֹא יָזוּן אֶת בְּהֶמְתּוֹ\hdot
בֵּין טְמֵאָה בֵּין טְהוֹרָה\hdot
רִבִּי אֱלִיעֶזֶר אוֹמֵר זָן הַטְּמֵאָה\hdot
וְאֵינוֹ זָן  אֶת הַטְּהוֹרָה\hdot
אָֽמְרוּ לוֹ מַה בֵּין טְמֵאָה לַטְּהוֹרָה\hdot
אָמַר לָהֶן שֶׁהַטְּהוֹרָה נַפְשָׁהּ לַשָּׁמַיִם וְגוּפָהּ שֶׁלּוֹ\hdot
וּטְמֵאָה נַפְשָׁהּ וְגוּפָהּ לַשָּׁמַיִם\hdot
אָֽמְרוּ לוֹ אַף הַטְּמֵאָה נַפְשָׁהּ לַשָּׁמַיִם וְגוּפָהּ  שֶׁלּוֹ\hdot
שֶׁאִם יִרְצֶה הֲרֵי הוּא מוֹכְרָהּ לַגּוֹיִם אוֹ מַאֲכִילָהּ לַכְּלָבִים׃}

\halacha{הַמֻּדָּר הֲנָאָה מֵחֲבֵירוֹ וְנִכְנַס לְבַקְּרוֹ\hdot
עוֹמֵד אֲבָל לֹא יוֹשֵׁב\hdot
וּמְרַפְּאֵהוּ רְפוּאַת נֶפֶשׁ\hdot
אֲבָל לֹא רְפוּאַת מָמוֹן\hdot
וְרוֹחֵץ עִמּוֹ בְאַמְבָּטִי גְדוֹלָה אֲבָל לֹא בַקְּטַנָּה\hdot
וְיָשֵׁן עִמּוֹ בַמִּטָּה\hdot
רִבִּי יְהוּדָה אוֹמֵר בִּימוֹת הַחַמָּה אֲבָל לֹא בִימוֹת הַגְּשָׁמִים\hdot
מִפְּנֵי שֶׁהוּא מְהַנֵּהוּ וּמֵסֵב עִמּוֹ עַל הַמִּטָּה\hdot
וְאוֹכֵל עִמּו עַל הַשֻׁלְחָן אֲבָל לֹא מִן הַתַּמְחוּי\hdot
אֲבָל אוֹכֵל הוּא מִן הַתַּמְחוּי הַחוֹזֵר\hdot
לֹא יֹאכַל עִמּוֹ מִן הָאֵבוּס שֶׁלִּפְנֵי הַפּוֹעֲלִים\hdot
וְלֹא יַעֲשֶׁה עִמּוֹ בָּאֹמֶן\hdot
דִּבְרֵי רִבִּי מֵאִיר\hdot
וַחֲכָמִים אוֹמְרִים עוֹשֶׂה הוּא בְרָחוֹק מִמֶּנּוּ׃}

\halacha{הַמֻּדָּר הֲנָאָה מֵחֲבֵירוֹ לִפְנֵי שְׁבִיעִית לֹא יוֹרֵד לְתוֹךְ שָׂדֵהוּ\hdot
וְאֵינוֹ אוֹכֵל מִן הַנּוֹטוֹת\hdot
וּבַשְּׁבִיעִית אֵינוֹ יוֹרֵד לְתוֹךְ שָׂדֵהוּ\hdot
אֲבָל אוֹכֵל הוּא מִן הַנּוֹטוֹת\hdot
נָדַר הֵימֶנּוּ מַאֲכָל לִפְנֵי שְׁבִיעִית\hdot
יוֹרֵד לְתוֹךְ  שָׂדֵהוּ וְאֵינוֹ אוֹכֵל מִן הַפֵּירוֹת\hdot
וּבַשְּׁבִיעִית יוֹרֵד וְאוֹכֵל׃}

\halacha{הַמֻּדָּר הֲנָאָה מֵחֲבֵירוֹ\hdot
לֹא יַשְׁאִילֶנּוּ\hdot
וְלֹא יִשְׁאַל מִמֶּנּוּ\hdot
לֹא יַלְוֶנּוּ וְלֹא יַלְוֶה מִמֶּנּוּ\hdot
וְלֹא יִמְכּוֹר לוֹ\hdot
וְלֹא יִקַּח מִמֶּנּוּ\hdot
אָמַר לוֹ הַשְׁאִילֵנִי פָּרָתָךְ\hdot
אָמַר לוֹ אֵינָהּ פְּנוּיָה\hdot
אָמַר קוּנָם שָׁדִי שֶׁאֵינִי חוֹרֵשׁ בָּהּ לְעוֹלָם\hdot
אִם הָיָה דַּרְכּוֹ לַחֲרוֹשׁ\hdot
הוּא אָסוּר וְכָל אָדָם מוּתָרִין\hdot
וְאִם אֵין דַּרְכּוֹ לַחֲרוֹשׁ\hdot
הוּא וְכָל אָדָם אֲסוּרִין׃}

\halacha{הַמֻּדָּר הֲנָאָה מֵחֲבֵירוֹ\hdot
וְאֵין לוֹ מַה יֹּאכַל\hdot
הוֹלֵךְ אֵצֶל הַחַנְוָנִי וְאוֹמֵר\hdot
אִישׁ פְּלוֹנִי מֻדָּר מִמֶּנִּי הֲנָאָה\hdot
וְאֵינִי יוֹדֵעַ מָה אֶעֱשֶׂה\hdot
הוּא נוֹתֵן לוֹ\hdot
וּבָא וְנוֹטֵל מִזֶּה\hdot
הָיָה בֵיתוֹ לִבְנוֹת\hdot
גִּדְרוֹ לִגְדּוֹר\hdot
שָׁדֵהוּ לִקְצוֹר\hdot
הוֹלֵךְ אֵצֶל הַפּוֹעֲלִים וְאוֹמֵר\hdot
אִישׁ פְּלוֹנִי מֻדָּר מִמֶּנִּי הֲנָאָה\hdot
וְאֵינִי יוֹדֵעַ מָה אֶעֱשֶׂה\hdot
הֵם עוֹשִׂים עִמּוֹ\hdot
וּבָאִין וְנוֹטְלִין שָׂכָר מִזֶּה׃}

\halacha{הָיוּ מְהַלְּכִין בַּדֶּרֶךְ וְאֵין לוֹ מַה יֹּאכַל\hdot
נוֹתֵן לְאַחֵר לְשֵׁם מַתָּנָה\hdot
וְהַלָּה מוּתָר בָּהּ\hdot
אִם אֵין עִמָּהֶם אַחֵר\hdot
מַנִּיחַ עַל הַסֶּלַע אוֹ עַל הַגָּדֵר וְאוֹמֵר הֲרֵי הֵן מֻפְקָרִין לְכָל מִי שֶׁיַּחְפּוֹץ\hdot
וְהַלָּה נוֹטֵל וְאוֹכֵל\hdot
וְרִבִּי יוֹסֵי אוֹסֵר׃}

\perek{הַשּׁוּתָפִין}{שֶׁנָּֽדְרוּ הֲנָאָה זֶה מִזֶּה\hdot
אֲסוּרִין לִכָּנֵס לֶחָצֵר\hdot
רִבִּי אֱלִיעֶזֶר בֶּן יַעֲקֹב אוֹמֵר זֶה נִכְנָס לְתוֹךְ שֶׁלּוֹ\hdot
וְזֶה נִכְנָס לְתוֹךְ שֶׁלּוֹ\hdot
וּשְׁנֵיהֶם אֲסוּרִים לְהַעֲמִיד שָׁם רֵיחַיִם וְתַנּוּר\hdot
וּלְגַדֵּל תַּרְנְגוֹלִים\hdot
הָיָה אֶחָד מֵהֶם מֻדָּר הֲנָאָה מֵחֲבֵירוֹ\hdot
לֹא יִכָּנֵס לֶחָצֵר\hdot
רִבִּי אֱלִיעֶזֶר בֶּן יַעֲקֹב אוֹמֵר יָכוֹל הוּא לוֹמַר לוֹ\hdot
לְתוֹךְ שֶׁלִּי אֲנִי נִכְנָס\hdot
וְאֵינִי נִכְנָס לְתוֹךְ שֶׁלָּךְ\hdot
וְכוֹפִין אֶת הַנּוֹדֵר לִמְכּוֹר אֶת חֶלְקוֹ׃}

\halacha{הָיָה אֶחָד מִן הַשּׁוּק מֻדָּר מֵאֶחָד מֵהֶם הֲנָאָה\hdot
לֹא יִכָּנֵס לֶחָצֵר\hdot
רִבִּי אֱלִיעֶזֶר בֶּן יַעֲקֹב אוֹמֵר יָכוֹל הוּא לוֹמַר לוֹ\hdot
לְתוֹךְ שֶׁל חֲבֵרָךְ אֲנִי נִכְנָס\hdot
וְאֵינִי נִכְנָס לְתוֹךְ שֶׁלָּךְ׃}

\halacha{הַמֻּדָּר הֲנָאָה מֵחֲבֵירוֹ\hdot
וְיֵשׁ לוֹ מֶרְחָץ וּבֵית הַבָּד מֻשְׂכָּרִים בָּעִיר\hdot
אִם יֵשׁ לוֹ בָהֶן תְּפִיסַת יָד אָסוּר\hdot
אֵין לוֹ  בָהֶן תְּפִיסַת יָד מוּתָר\hdot
הָאוֹמֵר לַחֲבֵירוֹ קוּנָם לְבֵיתְךָ שֶׁאֵינִי נִכְנָס\hdot
וְשָֽׁדְךָ שֶׁאֵינִי לוֹקֵחַ\hdot
מֵת אוֹ שֶׁמָּֽכְרָה לְאַחֵר מוּתָר\hdot
קוּנַם לְבַיִת זֶה שֶׁאֵינִי נִכְנָס\hdot
שָׁדֶה זוֹ שֶׁאֵינִי לוֹקֵחַ\hdot
מֵת אוֹ מְכָרָהּ לְאַחֵר\hdot
אָסוּר׃}

\halacha{הֲרֵינִי עָלֶיךָ חֵרֶם\hdot
הַמֻּדָּר אָסוּר\hdot
הֲרֵי אַתְּ עָלַי חֵרֶם הַנּוֹדֵר אָסוּר\hdot
הֲרֵינִי עָלֶיךָ וְאַתְּ עָלַי\hdot
שְׁנֵיהֶם אֲסוּרִין\hdot
וּשְׁנֵיהֶם מוּתָרִים בְּדָבָר שֶׁל עוֹלֵי בָבֶל\hdot
וַאֲסוּרִים בְּדָבָר שֶׁל אוֹתָהּ הָעִיר׃}

\halacha{וְאֵיזֶהוּ דָבָר שֶׁל עוֹלֵי בָבֶל\hdot
כְּגוֹן הַר הַבַּיִת וְהָעֲזָרוֹת וְהַבּוֹר שֶׁבְּאֶמְצַע הַדֶּרֶךְ\hdot
וְאֵיזֶהוּ דָבָר שֶׁל אוֹתָהּ הָעִיר\hdot
כְּגוֹן הָֽרְחוֹב וְהַמֶּרְחָץ וּבֵית הַכְּנֶסֶת\hdot
וְהַתֵּבָה וְהַסְּפָרִים\hdot
וְהַכּוֹתֵב חֶלְקוֹ לַנָּשִׂיא\hdot
רִבִּי יְהוּדָה אוֹמֵר אֶחָד כּוֹתֵב לַנָּשִׂיא\hdot
וְאֶחָד כּוֹתֵב לַהֶדְיוֹט\hdot
מַה בֵּין כּוֹתֵב לַנָּשִׂיא לַכּוֹתֵב לַהֶדְיוֹט\hdot
שֶׁהַכּוֹתֵב לַנָּשִׂיא אֵינוֹ צָרְיךְ לְזַכּוֹת\hdot
וְהַכּוֹתֵב לַהֶדְיוֹט צָרִיךְ לְזַכּוֹת\hdot
וַחֲכָמִים אוֹמְרִים אֶחָד זֶה וְאֶחָד זֶה צְרִיכִין לְזַכּוֹת\hdot
לֹא דִבְּרוּ בַנָּשִׂיא אֶלָּא בַהוֹוֶה\hdot
רִבִּי יְהוּדָה אוֹמֵר אֵין אַנְשֵׁי הַגָּלִיל צְרִיכִין לִכְתּוֹב שֶׁכְּבָר כָּֽתְבוּ אֲבוֹתֵיהֶם עַל יְדֵיהֶם׃}

\halacha{הַמֻּדָּר הֲנָאָה מֵחֲבֵירוֹ\hdot
וְאֵין לוֹ מַה יֹּאכַל\hdot
נוֹתֵן לְאַחֵר לְשׁוּם מַתָּנָה\hdot
וְהַלָּה מוּתָר בָּהּ\hdot
מַעֲשֶׂה בְאֶחָד בְּבֵית חוֹרוֹן שֶׁהָיָה אָבִיו מֻדָּר הֵימֶנּוּ הֲנָאָה\hdot
וְהָיָה מַשִּׂיא אֶת בְּנוֹ\hdot
וְאָמַר לַחֲבֵירוֹ חָצֵר וּסְעודָה נְתוּנִים לָךְ בְּמַתָּנָה\hdot
וְאֵינָן לְפָנֶיךָ אֶלָּא כְּדֵי שֶׁיָּבֹא אַבָּא\hdot
וְיֹאכַל עִמָּנוּ בַּסְּעוּדָה\hdot
אָמַר לוֹ אִם שֶׁלִּי הֵם\hdot
הֲרֵי הֵם מֻקְדָּשִׁים לַשָּמַיִם\hdot
אָמַר (לוֹ) לֹא נָתַתִּי לָךְ אֶת שֶׁלִּי\hdot
שֶׁתַּקְדִּישֵׁם לַשָּׁמַיִם\hdot
אָמַר לוֹ לֹא נָתַתָּ לִי אֶת שֶׁלָּךְ\hdot
אֶלָּא שְׁתְּהֵא אַתָּה וְאָבִיךָ אוֹכְלִים וְשׁוֹתִים וּמִתְרַצִּים זֶה לָזֶה\hdot
וִיהֵא עָוֺן תָּלוּי בְּרֹאשׁוֹ\hdot
(וּכְשֶׁבָּא דָבָר לִפְנֵי חֲכָמִים אָֽמְרוּ)
כָּל מַתָּנָה שֶׁאֵינָהּ שֶׁאִם הִקְדִישָׁהּ תְּהֵא מֻקְדֶּשֶׁת אֵינָהּ מַתָּנָה׃}

\perek{הַנּוֹדֵר}{מִן הַמְּבֻשָּׁל מוּתָר בִּצְלִי וּבְשָׁלוּק\hdot
אָמַר קוּנָם תַּבְשִׁיל שֶׁאֵינִי טוֹעֵם\hdot
אָסוּר בְּמַעֲשֵׂה קְדֵרָה רַךְ\hdot
וּמוּתָר בֶּעָבָה\hdot
וּמוּתָר בְּבֵיצַת טַרְמִיטָה\hdot
וּבְדָלַעַת הָֽרְמוּצָה׃}

\halacha{הַנּוֹדֵר מִמַּעֲשֵׂה קְדֵרָה אֵינוֹ אָסוּר אֶלָּא מִמַּעֲשֵׂה רְתִיחָתָהּ\hdot
אָמַר קוּנָם הַיּוֹרֵד לַקְּדֵרָה שֶׁאֵינִי טוֹעֵם\hdot
אָסוּר בְּכָל הַמִּתְבַּשְּׁלִין בַּקְּדֵרָה׃}

\halacha{מִן הַכִּבּוּשׁ\hdot
אֵינוֹ אָסוּר אֶלָּא מִן הַכִּבּוּשׁ שֶׁל יָרָק\hdot
כִּבּוּשׁ שֶׁאֵינִי טוֹעֵם אָסוּר בְּכָל הַכִּבּוּשִׁים\hdot
מִן הַשָּׁלוּק\hdot
אֵינוֹ אָסוּר אֶלָא מִן הַשָּלוּק שֶׁל בָּשָׂר\hdot
שָׁלוּק שֶׁאֵינִי טוֹעֵם\hdot
אָסוּר בְּכָל הַשְּׁלוּקִים\hdot
מִן הַצָּלִי\hdot
אֵינוֹ אָסוּר אֶלָּא מִן הַצָּלִי שֶׁל בָּשָׂר (דִּבְרֵי רִבִּי יְהוּדָה)\hdot
צָלִי שֶׁאֵינִי טוֹעֵם אָסוּר בְּכָל הַצְּלוּיִּן\hdot
מִן הַמָּלִיחַ\hdot
אֵינוֹ אָסוּר אֶלָּא מִן הַמָּלִיחַ שֶׁל דָּג\hdot
מָלִיחַ שֶׁאֵינִי טוֹעֵם\hdot
אָסוּר בְּכָל הַמְּלוּחִים׃}

\halacha{דָּג דָּגִים שֶׁאֵינִי טוֹעֵם\hdot
אָסוּר בָּהֶן בֵּין גְּדוֹלִים בֵּין קְטַנִּים\hdot
בֵּין מְלוּחִים בֵּין טְפֵלִים\hdot
בֵּין חַיִּין בֵּין מְבֻשָּׁלִין\hdot
וּמוּתָר בְּטָרִית טְרוּפָה וּבַצִּיר\hdot
הַנּוֹדֵר מִן הַצַּחֲנָה אָסוּר בְּטָרִית טְרוּפָה\hdot
וּמוּתָר בַּצִּיר וּבַמּוּרְיָס\hdot
הַנּוֹדֵר מִטָּרִית טְרוּפָה אָסוּר בַּצִּיר וּבַמּוּרְיָס׃}

\halacha{הַנּוֹדֵר מִן הֶחָלָב מוּתָר בַּקּוֹם\hdot
וְרִבִּי יוֹסֵי אוֹסֵר\hdot
מִן הַקּוֹם מוּתָר בֶּחָלָב\hdot
אַבָּא שָׁאוּל אוֹמֵר הַנּוֹדֵר מִן הַגְּבִינָה\hdot
אָסוּר בָּהּ בֵּין מְלוּחָה בֵּין טְפֵילָה׃}

\halacha{הַנּוֹדֵר מִן הַבָּשָׂר\hdot
מוּתָר בְּרֹטֶב וּבַקִּפָּה\hdot
וְרִבִּי יְהוּדָה אוֹסֵר\hdot
אָמַר רִבִּי יְהוּדָה מַעֲשֶׂה וְאָסַר (עָלַי) רִבִּי טַרְפוֹן בֵּצִים שֶׁנִּתְבַּשְּׁלוּ עִמּוֹ\hdot
אָֽמְרוּ לוֹ וְכֵן הַדָּבָר\hdot
אֵימָתַי\hdot
בִּזְמָן שֶׁיֹּאמַר בָּשָׂר זֶה עָלַי\hdot
שֶׁהַנּוֹדֵר מִן הַדָּבָר וְנִתְעָרַב בְּאַחֵר אִם יֵשׁ בּוֹ בְּנוֹתֵן טַעַם אָסוּר׃}

\halacha{הַנּוֹדֵר מִן הַיַּיִן\hdot
מוּתָר בְּתַבְשִׁיל שֶׁיֵּשׁ בּוֹ טַעַם יַיִן\hdot
אָמַר קוּנָם יַיִן זֶה שֶׁאֵינִי טוֹעֵם\hdot
וְנָפַל לַתַּבְשִׁיל\hdot
אִם יֵשׁ בּוֹ בְּנוֹתֵן טַעַם הֲרֵי זֶה אָסוּר\hdot
הַנּוֹדֵר מִן הָעֲנָבִים מוּתָר בַּיַּיִן\hdot
מִן הַזֵּיתִים מוּתָר בַּשֶּׁמֶן\hdot
אָמַר קוּנָם זֵיתִים וַעֲנָבִים אֵלּוּ שֶׁאֵינִי טוֹעֵם אָסוּר בָּהֶן וּבַיּוֹצֵא מֵהֶן׃}

\halacha{הַנּוֹדֵר מִן הַתְּמָרִים\hdot
מוּתָר בִּדְבַשׁ תְּמָרִים\hdot
מִסְּתַוְנִיּוֹת מוּתָר בְּחֹמֶץ סְתַוְנִיּות\hdot
רִבִּי יְהוּדָה בֶן בְּתֵירָה אוֹמֵר כָּל שֶׁשֵּׁם תּוֹלְדָתוֹ קְרוּיָה עָלָיו\hdot
וְנָדַר הֵימֶנּוּ\hdot
אָסוּר אַף בַּיּוֹצֵא הֵימֶנּוּ\hdot
וַחֲכָמִים מַתִּירִין׃}

\halacha{הַנּוֹדֵר מִן הַיַּיִן\hdot
מוּתָר בְּיֵין תַּפּוּחִים\hdot
מִן הַשֶּׁמֶן מוּתָר בְּשֶׁמֶן שׁוּמְשְׁמִין\hdot
מִן הַדְּבַשׁ\hdot
מוּתָר בִּדְבַשׁ תְּמָרִים\hdot
מִן הַחֹמֶץ\hdot
מוּתָר בְּחֹמֶץ סְתַוְנִיּוֹת\hdot
מִן הַכָּרִישִׁין מוּתָר בַּקַּפְלוֹטוֹת\hdot
מִן הַיָּרָק מוּתָר בִּירָקוֹת הַשָּׂדֶה\hdot
מִפְּנֵי שֶׁהוּא שֵׁם לִוּוּי׃}

\halacha{הַנּוֹדֵר מִן הַכְּרוּב אָסוּר בָּאֶסְפַּרְגּוֹס\hdot
מִן הָאֶסְפַּרְגּוֹס מוּתָר בַּכְּרוּב\hdot
מִן הַגְּרִיסִין אָסוּר בַּמַּקְפָּה\hdot
וְרִבִּי יוֹסֵי מַתִּיר\hdot
מִן הַמַּקְפָּה\hdot
מוּתָר בַּגְּרִיסִין\hdot
מִן הַמַּקְפָּה אָסוּר בְּשׁוּם\hdot
וְרִבִּי יוֹסֵי מַתִּיר\hdot
מִן הַשּׁוּם מוּתָר בַּמַּקְפָּה\hdot
מִן הָעֲדָשִׁים אָסוּר בָּאֲשִׁישִׁים\hdot
וְרִבִּי יוֹסֵי מַתִּיר\hdot
מִן הָאֲשִׁישִׁים מוּתָר בָּעֲדָשִׁים\hdot
חִטָּה (חִטִּים) שֶׁאֵינִי טוֹעֵם\hdot
אָסוּר בָּהֶן בֵּין קֶמַח בֵּין פַּת\hdot
גְּרִיס גְּרִיסִין שֶׁאֵינִי טוֹעֵם\hdot
אָסוּר בָּהֶן בֵּין חַיִּין בֵּין מְבֻשָּׁלִין\hdot
רִבִּי יְהוּדָה אוֹמֵר קוּנָם גְּרִיס אוֹ חִטִּים שֶׁאֵינִי טוֹעֵם מוּתָר לָכוֹס חַיִּין׃}

\perek{הַנּוֹדֵר}{מִן הַיָּרָק מוּתָר בַּדְּלוֹעִין וְרִבִּי עֲקִיבָא אוֹסֵר\hdot
אָֽמְרוּ לוֹ לְרִבִּי עֲקִיבָא\hdot
וַהֲלֹא אוֹמֵר אָדָם לִשְׁלוּחוֹ\hdot
קַח לִי יָרָק\hdot
וְהוּא אוֹמֵר לֹא מָצָאתִי אֶלָּא דְלוֹעִין אָמַר לָהֶם כֵּן הַדָּבָר\hdot
אוֹ שֶׁמֶּא אוֹמֵר הוּא לֹא מָצָאתִי אֶלָּא קָטְנִית\hdot
אֶלָּא שֶׁהַדְּלוֹעִים בִּכְלַל יָרָק (וְקָטְנִית אֵינוֹ בִּכְלַל יָרָק\hdot)
וְאָסוּר בְּפוֹל הַמִּצְרִי לַח\hdot
וּמוּתָר בַּיָּבֵשׁ׃}

\halacha{הַנּוֹדֵר מִן הַדָּגָן אָסוּר בְּפוֹל הַמִּצְרִי יָבֵשׁ דִּבְרֵי רִבִּי מֵאִיר\hdot
וַחֲכָמִים אוֹמְרִים אֵינוֹ אָסוּר אֶלָּא מֵחֲמֵשֶׁת הַמִּינִין\hdot
רִבִּי מֵאִיר אוֹמֵר הַנּוֹדֵר מִן הַתְּבוּאָה אֵינוֹ אָסוּר אֶלָּא מֶחֲמֶשֶׁת הַמִּינִין\hdot
אֲבָל (הַנּוֹדֵר) מִן הַדָּגָן אָסוּר בַּכֹּל וּמוּתָר בְּפֵירוֹת הָאִילָן וּבַיָּרָק׃}

\halacha{הַנּוֹדֵר מִן הַכְּסוּת\hdot
מוּתָר בַּשַּׂק\hdot
וּבַיְּרִיעָה וּבַחֲמִילָה\hdot
אָמַר קוּנָם צֶמֶר עוֹלֶה עָלַי\hdot
מוּתָר לְהִתְכַסּוֹת בְּגִיזֵי צֶמֶר\hdot
פִּשְׁתָּן עוֹלֶה עָלַי מוּתָר לְהִתְכַּסּוֹת בַּאֲנִיצֵי פִשְׁתָּן\hdot
רִבִּי יְהוּדָה אוֹמֵר הַכֹּל לְפִי הַנּוֹדֵר\hdot
טָעַן וְהִזִּיעַ וְהָיָה רֵיחוֹ קָשֶׁה\hdot
אָמַר קוּנָם צֶמֶר וּפִשְׁתִּים עוֹלִים עָלַי\hdot
מוּתָר לְהִתְכַּסּוֹת וְאָסוּר לְהַפְשִׁיל לַאֲחוֹרָיו׃}

\halacha{הַנּוֹדֵר מִן הַבַּיִת\hdot
מוּתָר בָּעֲלִיָּה\hdot
דִּבְרֵי רִבִּי מֵאִיר\hdot
וַחֲכָמִים אוֹמְרִים עֲלִיָּה בִּכְלַל הַבַּיִת\hdot
הַנּוֹדֵר מִן הָעֲלִיָּה מוּתָר בַּבַּיִת׃}

\halacha{הַנּוֹדֵר מִן הַמִּטָּה מוּתָר בַּדַּרְגָּשׁ\hdot
דִּבְרֵי רִבִּי מֵאִיר\hdot
וַחֲכָמִים אוֹמרִים דַּרְגָּש בִּכְלַל מִטָּה\hdot
הַנּוֹדֵר מִדַּרְגָּשׁ מוּתָר בַּמִּטָּה\hdot
הַנּוֹדֵר מִן הָעִיר מוּתָר לִכָּנֵס לִתְחוּמָהּ שֶׁל עִיר\hdot
וְאָסוּר לִכָּנֵס לְעִבּוּרָהּ\hdot
אֲבָל הַנּוֹדֵר מִן הַבַּיִת\hdot
אָסוּר מִן הָאָגָף וְלִפְנִים׃}

\halacha{קוּנָם פֵּירוֹת הָאֵלּוּ עָלַי\hdot
קוּנָם הֵן עַל פִּי\hdot
קוּנָם הֵן לְפִי\hdot
אָסוּר בְּחִלּוּפֵיהֶן וּבְגִדּוּלֵיהֶן\hdot
שֶׁאֵינִי אוֹכֵל וְשֶׁאֵינִי טוֹעֵם\hdot
מוּתָר בְּחִלּוּפֵיהֶן וּבְגִדּוּלֵיהֶן\hdot
בְּדָבָר שֶׁזַּרְעוֹ כָלֶה\hdot
אֲבָל בְּדָבָר שֶׁאֵין זַרְעוֹ כָלֶה אֲפִילוּ גִדּוּלֵי גִדּוּלִין אֲסוּרִין׃}

\halacha{הָאוֹמֵר לְאִשְׁתּוֹ קוּנָם מַעֲשֵׂה יָדַיִךְ עָלַי\hdot
קוּנָם הֵן עַל פִּי\hdot
קוּנָם הֵן לְפִי\hdot
אָסוּר בְּחִלּוּפֵיהֶן וּבְגִדּוּלֵיהֶן\hdot
שֶׁאֵינִי אוֹכֵל\hdot
שֶׁאֵינִי טוֹעֵם\hdot
מוּתָר בְּחִלּוּפֵיהֶן וּבְגִדּוּלֵיהֶן\hdot
בְּדָבָר שֶׁזַּרְעוֹ כָלֶה\hdot
אֲבָל בְּדָבָר שֶׁאֵין זַרְעוֹ כָלֶה אֲפִילוּ גִדּוּלֵי גְדּוּלִין אֲסורִין}

\halacha{שֶׁאַתְּ עוֹשָׂה אֵינִי אוֹכֵל עַד הַפֶּסַח\hdot
שֶׁאַתְּ עוֹשָׂה אֵינִי מִתְכַסֶּה עַד הַפֶּסַח\hdot
עָֽשְׂתָה לִפְנֵי הַפֶּסַח מוּתָר לֶאֱכוֹל וּלְהִתְכַּסּוֹת אַחַר הַפֶּסַח\hdot
שֶׁאַתְּ עוֹשָׂה עַד הַפֶּסַח אֵינִי אוֹכֵל\hdot
וּשְׁאַתְּ עוֹשָׂה עַד הַפֶּסַח אֵינִי מִתְכַּסֶּה\hdot
עָֽשְׂתָה לִפְנֵי הַפֶּסַח אָסוּר לֶאֱכוֹל וּלְהִתְכַּסּוֹת אַחַר הַפֶּסַח׃}

\halacha{שֶׁאַתְּ נֶהֱנֵית לִי עַד הַפֶּסַח אִם תֵּלְכִי לְבֵית אַבִיךְ עַד הֶחָג\hdot
הָֽלְכָה לִפְנֵי הַפֶּסַח\hdot
אֲסוּרָה בַהֲנָאָתוֹ עַד הַפֶּסַח\hdot
אַחַר הַפֶּסַח בְּלֹא יַחֵל דְּבָרוֹ\hdot
שֶׁאַתְּ נֶהֱנֵית לִי עַד הֶחָג אִם תֵּלְכִי לְבֵית אָבִיךְ עַד הַפֶּסַח\hdot
וְהָֽלְכָה לִפְנֵי הַפֶּסַח אֲסוּרָה בַהֲנָאָתוֹ עַד הֶחָג\hdot
וּמוּתֶרֶת לִילֵךְ אַחַר הַפֶּסַח׃}

\perek{קוּנָם}{יַיִן שֶׁאֵינִי טוֹעֶם הַיּוֹם\hdot
אֵינוֹ אָסוּר אֶלָּא עַד שֶׁתֶּחְשַׁךְ\hdot
שַׁבָּת זוֹ אָסוּר בְּכָל השַּׁבָּת\hdot
וְשַׁבָּת שֶׁעָֽבְרָה\hdot
חֹדֶשׁ זֶה\hdot
אָסוּר בְּכָל הַחֹדֶשׁ\hdot
וְרֹאשׁ חֹדֶשׁ לְהַבָּא\hdot
שָׁנָה זוֹ\hdot
אָסוּר בְּכָל הַשָּׁנָה\hdot
וְרֹאשׁ הַשָּׁנָה לֶעָתִיד לָבֹא\hdot
שָׁבוּעַ זֶה\hdot
אָסוּר בְּכָל הַשָּׁבוּעַ\hdot
וּשְׁבִיעִית שֶׁעָֽבְרָה\hdot
וְאִם אָמַר יוֹם אֶחָד\hdot
שַׁבָּת אֶחָד\hdot
חֹדשֶ אֶחָד\hdot
שָׁנָה אֶחָת\hdot
שָׁבוּעַ אֶחָד\hdot
אָסוּר מִיּוֹם לְיוֹם׃}

\halacha{עַד הַפֶּסַח\hdot
אָסוּר עַד שֶׁיַּגִּיעַ\hdot
עַד שֶׁיְּהֵא\hdot
אָסוּר עַד שֶׁיֵּצֵא\hdot
עַד לִפְנֵי הַפֶּסַח\hdot
רִבִּי מֵאִיר אוֹמֵר אָסוּר עַד שֶׁיַּגִּיעַ\hdot
רִבִּי יוֹסֵי אוֹמֵר עַד שֶׁיֵּצֵא׃}

\halacha{עַד הַקָּצִיר\hdot
עַד הַבָּצִיר\hdot
עַד הַמָּסִיק\hdot
אֵינוֹ אָסוּר אֶלָּא עַד שֶׁיַּגִּיעַ\hdot
זֶה הַכְּלָל\hdot
כָּל שֶׁזְּמַנּוֹ קָבוּעַ וְאָמַר עַד שֶׁיַּגִּיעַ אָסוּר ֹעַד שֶׁיַּגִּיעַ\hdot
אָמַר עַד שְׁיְּהֵא אָסוּר עַד שֶׁיֵּצֵא\hdot
וְכָל שֶׁאֵין זְמַנּוֹ קָבוּעַ\hdot
בֵּין אָמַר עַד שֶׁיְּהֵא בֵּין אָמַר עַד שֶׁיַּגִּיעַ אֵינוֹ אָסוּר אֶלָּא עַד שֶׁיַּגִּיעַ׃}

\halacha{עַד הַקַּיִץ\hdot
עַד שֶׁיְּהֵא הַקַּיִץ\hdot
עַד שֶׁיַּתְחִילוּ הָעָם לְהַכְנִיס בַּכַּלְכָּלוֹת\hdot
עַד שֶׁיַּעֲבוֹר הַקַּיִץ\hdot
עַד שֶׁיִּקְפְּלוּ הַמִּקְצוֹעוֹת\hdot
עַד הַקָּצִיר עַד שֶׁיַּתְחִיל הָעָם לִקְצוֹר קְצִיר חִטִּים אֲבָל לֹא קְצִיר שְׂעוֹרִים\hdot
הַכֹּל לְפִי (מְקוֹם) נִדְרוֹ\hdot
אִם הָיָה בָהָר בָּהָר\hdot
וְאִם הָיָה בַבִּקְעָה בַּבִּקְעָה׃}

\halacha{עַד הַגְּשָׁמִים\hdot
עַד שֶׁיִּהְיוּ הַגְּשָׁמִים\hdot
עַד שֶׁתֵּרֵד רְבִיעָה שְׁנִיָּה\hdot
רַבָּן שִׁמְעוׁן בֶּן גַּמְלִיאֵל אוֹמֵר עַד שֶׁיַּגִּיעַ זְמַנָּהּ שֶׁל רְבִיעָה\hdot
עַד שֶׁיִּפְסְקוּ גְשָׁמִים\hdot
עַד שֶׁיֵּצֵא נִיסָן כֻּלּוֹ\hdot
דִּבְרֵי רִבִּי מֵאִיר\hdot
רִבִּי יְהוּדָה אוֹמֵר עַד שֶׁיַּעֲבוֹר הַפֶּסַח\hdot
קוּנָם יַיִן שֶׁאֵינִי טוֹעֵם הַשָּׁנָה\hdot
נִתְעַבְּרָה הַשָּׁנָה אָסוּר בָּהּ וּבְעִבּורָהּ\hdot
עַד רֹאשׁ אֲדָר\hdot
עַד רֹאשׁ אֲדָר הָרִאשׁוֹן\hdot
עַד סוֹף אֲדָר\hdot
עַד סוֹף אֲדָר רִאשׁוֹן\hdot
רִבִּי יְהוּדָה אוֹמֵר קוּנָם יַיִן שֶׁאֵינִי טוֹעֵם עַד שֶׁיְּהֵא הַפֶּסַח\hdot
אֵינוֹ אָסוּר אֶלָּא עַד לֵילֵי הַפֶּסַח\hdot
שֶׁלֹּא נִתְכַּוָּן זֶה אֶלָּא עַד שָׁעָה שֶׁדֶּרֶךְ בְּנֵי אָדָם לִשְׁתּוֹת יַיִן׃}

\halacha{אָמַר קוּנָם בָּשָׂר שֶׁאֵינִי טוֹעֵם עַד שֶׁיְּהֵא הַצּוֹם\hdot
אֵינוֹ אָסוּר אֶלָּא עַד לֵילֵי הַצּוֹם שֶׁלֹּא נִתְכַּוָּן זֶה אֶלָּא עַד שָׁעָה שֶׁדֶּרֶךְ בְּנֵי אָדָם לֶאֱכוֹל בָּשָׂר\hdot
רִבִּי יוֹסֵי בְּנוֹ אוֹמֵר אָמַר קוּנָם שׁוּם שֶׁאֵינִי טוֹעֵם\hdot
עַד שֶׁתְּהֵא שַׁבָּת\hdot
אֵינוֹ אָסוּר אֶלָּא עַד לֵילֵי שַׁבָּת\hdot
שֶׁלֹּא נִתְכַּוָּן זֶה אֶלָּא עַד שָׁעָה שֶׁדֶּרֶךְ בְּנֵי אָדָם לֶאֱכוֹל שׁוּם׃}

\halacha{הָאוֹמֵר לַחֲבֵירוֹ קוּנָם שֶׁאֵינִי נֶהֱנֶה לָךְ\hdot
אִם אֵין אַתָּה בָא וְנוֹטֵל לְבָנֶיךָ\hdot
כּוֹר אֶחָד שֶׁל חִטִּין וּשְׁתֵּי חֲבִיּוֹת שֶׁל יַיִן\hdot
הֲרֵי זֶה יָכוֹל לְהָפֵר אֶת נִדְרוֹ\hdot
שֶׁלֹּא עַל פִּי חָכָם\hdot
וְיֹאמַר לוֹ כְּלוּם אָמַרְתָּ אֶלָּא מִפְּנֵי כְבוֹדִי\hdot
זֶהוּ כְבוֹדִי\hdot
וְכֵן הָאוֹמֵר לַחֲבֵירוֹ קוּנָם שֶׁאַתָּה נֶהֱנֶה לִי אִם אֵין אַתָּה בָא וְנוֹתֵן לִבְנִי כּוֹר אֶחָד שֶׁל חִטִּין\hdot
וּשְׁתֵּי חֲבִיּוֹת שֶׁל יַיִן רִבִּי מֵאִיר אוֹמֵר אָסוּר עַד שֶׁיִּתֵּן\hdot
וַחֲכָמִים אוֹמְרִים אַף זֶה יָכוֹל לְהָפֵר אֶת נִדְרוֹ שֶׁלֹּא עַל פִּי חָכָם\hdot
וְיֹאמַר לוֹ הֲרֵי אֲנִי כְּאִלּוּ הִתְקַבַּלְתִּי\hdot
הָיוּ מְסָֽרְבִין בּוֹ לִשָּׂא בַת אֲחוֹתוֹ\hdot
וְאָמַר קוּנָם שֶׁהִיא נֶהֱנֵית לִי לְעוֹלָם\hdot
וְכֵן הַמְּגָרֵשׁ אֶת אִשְׁתּוֹ\hdot
וְאָמַר קוּנָם אִשְׁתִּי נֶהֱנֵית לִי לְעוֹלָם\hdot
הֲרֵי אֵלּוּ מוּתָרוֹת לִהֲנוֹת לוֹ\hdot
שֶׁלֹּא נִתְכַּוָּן זֶה אֶלָּא לְשׁוּם אִישׁוּת\hdot
הָיָה מְסָרֵב בַּחֲבֵירוֹ שֶׁיֹּאכַל אֶצְלוֹ\hdot
אָמַר קוּנָם לְבֵיתְךָ שֶׁאֵינִי נִכְנָס\hdot
טִפַּת צוֹנֵן שֶׁאֵינִי טוֹעֵם לָךְ\hdot
מוּתָר לִכָּנֵס לְבֵיתוֹ\hdot
וְלִשְׁתּוֹת מִמֶּנּוּ צוֹנֵן שֶׁלֹּא נִתְכַּוָּן זֶה\hdot
אֶלָּא לְשׁוּם אֲכִילָה וּשְׁתִיָּה׃}

\perek{רִבִּי}{אֱלִיעֶזֶר אוֹמֵר פּוֹתְחִין לְאָדָם בִּכְבוֹד אָבִיו וְאִמּוֹ\hdot
וַחֲכָמִים אוֹסְרִין\hdot
אָמַר רִבִּי צָדוֹק עַד שֶׁפּוֹתְחִין לוֹ בִּכְבוֹד אָבִיו וְאִמּוֹ\hdot
יִפְתְּחוּ לוֹ בִּכְבוֹד הַמָּקוֹם (אָֽמְרוּ לוֹ) אִם כֵּן אֵין נְדָרִים\hdot
מוֹדִים חֲכָמִים לְרִבִּי אֱלִיעֶזֶר בְּדָבָר שֶׁבֵּינוֹ לְבֵין אָבִיו וְאִמּוֹ\hdot
שֶׁפּוֹתְחִין לוֹ בִּכְבוֹד אָבִיו וְאִמּוֹ׃}
<<<<<<< HEAD

\halacha{וְעוֹד אָמַר רִבִּי אֱלִיעֶזֶר פּוֹתְחִים בַּנּוֹלָד\hdot
וַחֲכָמִים אוֹסְרִין\hdot
כֵּיצַד\hdot
אָמַר קוּנָם שֶׁאֵינִי נֶהֱנֶה לְאִיש פְּלוֹנִי וְנַעֲשָׂה סוֹפֵר\hdot
אוֹ שֶׁהָיָה מַשִּׂיא אֶת בְּנוֹ (בְּקָרוֹב) וְאָמַר אִלּוּ הָיִיתִי יוֹדֵעַ שֶׁהוּא נַעֲשָׂה סוֹפֵר\hdot
אוֹ שֶׁהָיָה מַשִּׂיא אֶת בְּנוֹ בְּקָרוֹב\hdot
לֹא הָיִיתִי נוֹדֵר\hdot
קוּנָם לְבַיִת זֶה שֶׁאֵינִי נִכְנָס\hdot
וְנַעֲשָׂה בֵּית הַכְּנֶסֶת אָמַר אִלּוּ הָיִיתִי יוֹדֵעַ שֶׁהוּא נַעֲשָׂה בֵּית הַכְּנֶסֶת\hdot
לֹא הָיִיתִי נוֹדֵר\hdot
רִבִּי אֱלִיעֶזֶר מַתִּיר\hdot
וַחֲכָמִים אוֹסְרִין׃}

\halacha{רִבִּי מֵאִיר אוֹמֵר יֵשׁ דְּבָרִים שֶׁהֵן כַּנּוֹלָד וְאֵינָן כַּנּוֹלָד וַחֲכָמִים מוֹדִים לוֹ\hdot
כֵּיצַד\hdot
אָמַר קוּנָם שֶׁאֵינִי נוֹשֵׂא אֶת פְּלוֹנִית\hdot
שֶׁאָבִיהָ רָע\hdot
אָֽמְרוּ לוֹ מֵת\hdot
אוֹ שֶׂעָשָׂה תְשׁוּבָה\hdot
קוּנָם לְבַיִת זֶה שֶׁאֵינִי נִכְנָס שֶׁהַכֶּלֶב רָע בְּתוֹכוֹ\hdot
אוֹ שֶׁהַנָּחָשׁ בְּתוֹכוֹ\hdot
אָֽמְרוּ לוֹ מֵת הַכֶּלֶב\hdot
אוֹ שֶׁנֶהֱרַג הַנָּחָשׁ\hdot
הֲרֵי הֵן כַּנּוֹלָד וְאֵינָן כַּנּוֹלָד וַחֲכָמִים מוֹדִים לוֹ׃}
=======
>>>>>>> refs/remotes/origin/main

\tractatehead{מסכת נזיר}{nazir}

\perek{כָּל}{כִּנּוּיֵי נְזִירוּת\hdot
כִּנְזִירוּת\hdot
הָאוֹמֵר אֲהֵא הֲרֵי זֶה נָזִיר\hdot
אוֹ אֲהֵא נָוֶה\hdot
נָזִיר\hdot
נָזִיק\hdot
נָזִיחַ\hdot
פָּזִיחַ\hdot
הֲרֵי זֶה נָזִיר\hdot
הֲרֵינִי כַזֶּה\hdot
הֲרֵינִי מְסַלְסֵל\hdot
הֲרֵינִי מְכַלְכֵּל\hdot
הֲרֵי עָלַי לִשְׁלוֹחַ פֶּרַע\hdot
הֲרֵי זֶה נָזִיר\hdot
הֲרֵי עָלַי צִפּוֹרִים\hdot
רִבִּי מֵאִיר אוֹמֵר נָזִיר\hdot
וַחֲכָמִים אוֹמְרִים אֵינוֹ נָזִיר׃}

\halacha{הֲרֵינִי נָזִיר מִן הַחַרְצַנִּים וּמִן הַזַּגִּים וּמִן הַתִּגְלַחַת\hdot
וּמִן הַטֻּמְאָה\hdot
הֲרֵי זֶה נָזִיר\hdot
וְכָל דִּקְדוּקֵי נְזִירוּת עָלָיו\hdot
הֲרֵינִי כְשִׁמְשׁוֹן\hdot
כְּבֶן מָנוֹחַ\hdot
כְּבַעַל דְּלִילָה\hdot
כְּמִי שֶׁעָקַר דַּלְתּוֹת עַזָּה\hdot
כְּמִי שֶׁנָּֽקְרוּ פְלִשְׁתִּים אֶת עֵינָיו\hdot
הֲרֵי זֶה נְזִיר שְׁמְשׁוֹן\hdot
מַה בֵּין נְזִיר עוֹלָם לִנְזִיר שִׁמְשׁוֹן\hdot
נְזִיר עוֹלָם הִכְבִּיד שְׂעָרוֹ מֵיקֵל בְּתַעַר\hdot
וּמֵבִיא שָׁלֹשׁ בְּהֵמוֹת\hdot
וְאִם נִטְמָא מֵבִיא קָרְבַּן טֻמְאָה\hdot
נְזִיר שִׁמְשׁוֹן הִכְבִּיד שְׂעָרוֹ אֵינוֹ מֵיקֵל\hdot
וְאִם נִטְמָא אֵינוֹ מֵבִיא קָרְבָּן טֻמְאָה׃}

\halacha{סְתָם נְזִירוּת שְׁלֹשִׁים יוֹם\hdot
אָמַר הֲרֵינִי נָזִיר אַחַת גְּדוֹלָה\hdot
הֲרֵינִי נָזִיר אַחַת קְטַנָּה\hdot
אֲפִילוּ מִכָּאן וְעַד סוֹף הָעוֹלָם\hdot
נָזִיר שְׁלֹשִׁים יוֹם\hdot
הֲרֵינִי נָזִיר וְיוֹם אֶחָד\hdot
הֲרֵינִי נָזִיר וְשָׁעָה אַחַת\hdot
הֲרֵינִי נָזִיר אַחַת וּמֶחְצָה\hdot
הֲרֵי זֶה נָזִיר שְׁתַּיִם\hdot
הֲרֵינִי נָזִיר שְׁלֹשִׁים יוֹם וְשָׁעָה אַחַת\hdot
נָזִיר שְׁלֹשִׁים וְאֶחָד יוֹם\hdot
שֶׁאֵין נוֹזְרִים לְשָׁעוֹת׃}

\halacha{הֲרֵינִי נָזִיר כִּשְׂעַר רֹאשִׁי\hdot
וְכַעֲפַר הָאָרֶץ\hdot
וּכְחוֹל הַיָּם\hdot
הֲרֵי זֶה נְזִיר עוֹלָם\hdot
וּמְגַלֵּחַ אַחַת לִשְׁלֹשִׁים יוֹם\hdot
רִבִּי אוֹמֵר אֵין זֶה מְגַלֵּחַ אַחַת לִשְׁלֹשִׁים יוֹם\hdot
וְאֵיזֶה שֶׁמְּגַלֵּחַ אַחַת לִשְׁלֹשִׁים יוֹם\hdot
הָאוֹמֵר הֲרֵי עָלַי נְזִירוּת כִּשְׂעַר רֹאשִׁי\hdot
וְכַעֲפַר הָאָרֶץ\hdot
וּכְחוֹל  הַיָּם׃}

\halacha{הֲרֵינִי נָזִיר מְלֹא הַבַּיִת\hdot
אוֹ מְלֹא הַקֻּפָּה\hdot
בּוֹדְקִין אוֹתוֹ\hdot
אִם אָמַר אַחַת גְּדוֹלַה נָזַרְתִּי\hdot
נָזִיר שְׁלֹשִׁים יוֹם\hdot
וְאִם אָמַר סְתָם נָזַרְתִּי\hdot
רוֹאִין אֶת הַקֻּפָּה כְּאִלּוּ הִיא מְלֵאָה חַרְדָּל\hdot
וְנָזִיר כָּל יָמָיו׃}

\halacha{הֲרֵינִי נָזִיר מִכָּאן עַד מָקוֹם פְּלוֹנִי\hdot
אוֹמְדִים כַּמָּה יָמִים מִכָּאן עַד מָקוֹם פְּלוֹנִי\hdot
אִם פָּחוּת מִשְּׁלֹשִׁים יוֹם נָזִיר שְׁלֹשִׁים יוֹם\hdot
וְאִם לָאוּ נָזִיר כְּמִנְיַן הַיָּמִים׃}

\halacha{הֲרֵינִי נָזִיר כְּמִנְיַן יְמוֹת הַחַמָּה\hdot
מוֹנֶה נְזִירוּת כְּמִנְיַן יְמוֹת הַחַמָּה\hdot
אָמַר רִבִּי יְהוּדָה מַעֲשֶׂה הָיָה\hdot
וְכֵיוַן שֶׁהִשְׁלִים מֵת׃}

\perek{הֲרֵינִי}{נָזִיר מִן הַגְּרוֹגְרוֹת\hdot
וּמִן הַדְּבֵילָה\hdot
בֵּית שַׁמָּאי אוֹמְרִים נָזִיר\hdot
וּבֵית הִלֵּל אוֹמְרִים אֵינוֹ נָזִיר\hdot
אָמַר רִבִּי יְהוּדָה אַף כְּשֶׁאָֽמְרוּ בֵית שַׁמָּאי לֹא אָֽמְרוּ\hdot
אֶלָּא בְּאוֹמֵר הֲרֵי הֵן עָלַי כְּקָרְבָּן׃}

\halacha{אָמַר\hdot
אָֽמְרָה פָרָה זוֹ הֲרֵינִי נְזִירָה אִם עוֹמֶדֶת (אָנִי)\hdot
אָמַר הַדֶּלֶת הַזֶּה\hdot
הֲרֵינִי נָזִיר אִם נִפְתַּח אָנִי\hdot
בֵּית שַׁמָּאי אוֹמְרִים נָזִיר\hdot
וּבֵית הִלֵּל אוֹמְרִים אֵינוֹ נָזִיר\hdot
אָמַר רִבִּי יְהוּדָה\hdot
אַף כְּשֶׁאָֽמְרוּ בֵית  שַׁמָּאי לֹא אָֽמְרוּ\hdot
אֶלָּא בְּאוֹמֵר הֲרֵי פָרָה זוֹ קָרְבָּן עָלַי\hdot
אִם עוֹמֶדֶת הִיא׃}

\halacha{מָֽזְגוּ לוֹ אֶת הַכּוֹס\hdot
וְאָמַר הֲרֵינִי נָזִיר מִמֶּנּוּ\hdot
הֲרֵי זֶה נָזִיר\hdot
מַעֲשֶׂה בְּאִשָּׁה אַחַת שֶׁהָֽיְתָה שִׁכּוֹרָה\hdot
וּמָֽזְגוּ לָהּ אֶת הַכּוֹס\hdot
וְאָֽמְרָה הֲרֵינִי נְזִירָה מִמֶּנּוּ\hdot
אָֽמְרוּ חֲכָמִים לֹא נִתְכַּוְּנָה זוֹ\hdot
אֶלָּא לוֹמַר הֲרֵי הוּא עָלַי קָרְבָּן׃}

\halacha{הֲרֵינִי נָזִיר\hdot
עַל מְנַת  שֶׁאֶהֵא שׁוֹתֶה יַיִן וּמִטַּמֵּא לַמֵּתִים\hdot
הֲרֵי זֶה נָזִיר\hdot
וְאָסוּר בְּכֻלָּן\hdot
יוֹדֵעַ אֲנִי שֶׁיֵּשׁ נְזִירוּת\hdot
אֲבָל אֵינִי יוֹדֵעַ שֶׁהַנָּזִיר אָסוּר בַּיַּיִן\hdot
הֲרֵי זֶה אָסוּר\hdot
וְרִבִּי שִׁמְעוֹן מַתִּיר\hdot
יוֹדֵעַ אֲנִי שֶׁהַנָּזִיר אָסוּר בַּיַּיִן\hdot
אֲבָל סָבוּר הָיִיתִי שֶׁחֲכָמִים מַתִּירִין לִי\hdot
מִפְּנֵי שֶׁאֵין אֲנִי יָכוֹל לִחְיוֹת אֶלָּא בַיַּיִן\hdot
אוֹ מִפְּנֵי שֶׁאֲנִי קוֹבֵר אֶת הַמֵּתִים\hdot
הֲרֵי זֶה מוּתָר\hdot
וְרִבִּי שִׁמְעוֹן אוֹסֵר׃}


\clearpage

\thispagestyle{empty}
\centerlastline
\small
\parskip 1pt
{\centering
  \scriptsize
   אחר השלמת כל מסכתא ומסכתא יאמר זה שלש פעמים\\ ויועיל לשכחה בעזרת השם יתברך\hdot
}

הַדְרָן עֲלָךְ מַסֶּכֶת {\scriptsize(פלוני)} וְהַדְרָךְ עֲלָן\hdot
דַּעְתָן עֲלָךְ מַסֶּכֶת {\scriptsize(פלוני)} וְדַעְתָּך עֲלָן\hdot
לָא נִתְנְשֵׁי מִינָךְ מַסֶּכֶת {\scriptsize(פלוני}) וְלָא תִתְנְשֵׁי מִינָן
לָא בְעַלְמָא הַדֵּין וְלָא בְעַלְמָא דְאָתֵי׃

\vspace*{1em}

{\scriptsize ואחר כל יאמר}

יְהִי רָצוֹן מִלְּפָנֶיךָ יְיָ אֱלֹהֵינוּ וֵאלֹהֵי אֲבוֹתֵינוּ שֶׁתְּהֵא תוֹרָֽתְךָ אוּמָנוּתֵנוּ בָּעוֹלָם הַזֶּה וּתְהֵא עִמָּנו לְעוֹלָם הַבָּא\hdot
חֲנַנְיָה בַּר פַּפָּא\hdot
רָמֵי בַּר פַּפָּא\hdot
נַחְמָן בַּר פַּפָּא\hdot
אָחַאי בַּר פַּפָּא\hdot
אַבָּא מָרִי בַּר פַּפָּא\hdot
רַפְרָם בַּר פַּפָּא\hdot
רָכִישׁ בַּר פַּפָּא\hdot
סֻרְחַב בַּר פַּפָּא\hdot
אָדָא בַּר פַּפָּא\hdot
דָרוֹ בַּר פַּפָּא׃

\vspace*{1em}

הַעֲרֵב נָא יְיָ אֱלֹהֵינוּ אֶת דִּבְרֵי תוֹרָֽתְךָ בְפִינוּ וּבְפִיפִיּוֹת עַמְּךָ בֵּית יִשְׂרָאֵל וְנִהְיֶה אֲנַחְנוּ כֻּלָּנוּ וְצֶאֱצָאֵינוּ וְצֶאֱצָאֵי עַמְּךָ בֵּית יִשְׂרָאֵל יִרְאֵי שְׁמֶךָ וְלוֹמְדֵי תוֹרָֽתֶךָ׃

מֵאֹיְבַי תְּחַכְּמֵנִי מִצְוֺתֶךָ כִּי לְעוֹלָם הִיא לִי׃

יְהִי לִבִּי תָמִים בְּחֻקֶּיךָ לְמַעַן לֹא אֵבוֹשׁ׃

לְעוֹלָם לֹא אֶשְׁכַּח פִּקּוּדֶיךָ כִּי בָם חִיִּיתָנִי׃

  בָּרוּךְ אַתָּה יְהוָה לַמְּדֵנִי חֻקֶּיךָ׃


%\def \sederCount{שלישי}
%\clearpage

\centerlastline
{
\small
\parskip 1pt
{\centering
  \scriptsize
   אחר השלמת כל מסכתא ומסכתא יאמר זה שלש פעמים\\ ויועיל לשכחה בעזרת השם יתברך\hdot
}

\begin{tp}{219}
הַדְרָן עֲלָךְ מַסֶּכֶת {\scriptsize(פלוני)} וְהַדְרָךְ עֲלָן\hdot
דַּעְתָן עֲלָךְ מַסֶּכֶת {\scriptsize(פלוני)} וְדַעְתָּך עֲלָן\hdot
לָא נִתְנְשֵׁי מִינָךְ מַסֶּכֶת {\scriptsize(פלוני}) וְלָא תִתְנְשֵׁי מִינָן
לָא בְעַלְמָא הַדֵּין וְלָא בְעַלְמָא דְאָתֵי׃
\end{tp}

{\scriptsize ואחר כל יאמר}

יְהִי רָצוֹן מִלְּפָנֶיךָ יְיָ אֱלֹהֵינוּ וֵאלֹהֵי אֲבוֹתֵינוּ שֶׁתְּהֵא תוֹרָֽתְךָ אוּמָנוּתֵנוּ בָּעוֹלָם הַזֶּה וּתְהֵא עִמָּנו לְעוּלָם הַבָּא\hdot
חֲנַנְיָה בַּר פַּפָּא\hdot
רָמֵי בַּר פַּפָּא\hdot
נַחְמָן בַּר פַּפָּא\hdot
אָחַאי בַּר פַּפָּא\hdot
אַבָּא מָרִי בַּר פַּפָּא\hdot
רַפְרָם בַּר פַּפָּא\hdot
רָכִישׁ בַּר פַּפָּא\hdot
סֻרְחַב בַּר פַּפָּא\hdot
אָדָא בַּר פַּפָּא\hdot
דָרוֹ בַּר פַּפָּא׃

\hbadness 64
\begin{tp}{2048}
הַעֲרֵב נָא יְיָ אֱלֹהֵינוּ אֶת דִּבְרֵי תוֹרָֽתְךָ בְפִינוּ וּבְפִיפִיּוֹת עַמְּךָ בֵּית יִשְׂרָאֵל וְנִהְיֶה אֲנַחְנוּ כֻּלָּנוּ וְצֶאֱצָאֵינוּ וְצֶאֱצָאֵי עַמְּךָ בֵּית יִשְׂרָאֵל יִרְאֵי שְׁמֶךָ וְלוֹמְדֵי תוֹרָֽתֶךָ\hdot
 מֵאֹיְבַי תְּחַכְּמֵנִי מִצְו‍ֹתֶךָ כִּי לְעוֹלָם הִיא לִי\hdot
 יְהִי לִבִּי תָמִים בְּחֻקֶּיךָ לְמַעַן לֹא אֵבוֹשׁ\hdot
 לְעוֹלָם לֹא אֶשְׁכַּח פִּקּוּדֶיךָ כִּי בָם חִיִּיתָנִי\hdot
  בָּרוּךְ אַתָּה יְהוָה לַמְּדֵנִי חֻקֶּיךָ׃
\end{tp}
  
    {\centering
      \larger חתימת חלק ראשון מהמשנה

    }

    לְהוֹדוֹת לְהַלֵּל לְשַׁבֵּחַ לְפָאֵר לְרוֹמֵם לְהַדֵּר וּלְנַצֵּחַ עַל כָּל דִּבְרֵי שִׁירוֹת וְתִשְׁבָּחוֹת דָּוִד בֶּן יִשַּׁי עַבְדְּךָ מְשִׁיחֶךָ׃ בַּעֲצָֽתְךָ תַנְחֵנִי וְאַחַר כָּבוֹד תִּקָּחֵנִי׃
    וְהוּא רַחוּם יְכַפֵּר עָו‍ֹן וְלֹא יַשְׁחִית וְהִרְבָּה לְהָשִׁיב אַפּוֹ וְלֹא יָעִיר כָּל חֲמָתוֹ׃
    אַשְׁרֵי הַגֶּבֶר אֲשֶׁר תְּיַסְּרֶנּוּ יָּהּ וּמִתּוֹרָֽתְךָ תְלַמְּדֶנּוּ׃
    וַאֲנִי בְּחַסְדְּךָ בָטַחְתִּי יָגֵל לִבִּי בִּישׁוּעָתֶךָ אָשִׁירָה לַיְיָ כִּי גָמַל עָלָי׃
    שׂוֹשׂ אָשִׂישׂ בַּיְיָ תָּגֵל נַפְשִׁי בֵּאלֹהַי כִּי הִלְבִּישַׁנִי בִּגְדֵי יֶשַׁע מְעִיל צְדָקָה יְעָטָנִי כֶּחָתָן יְכַהֵן פְּאֵר וְכַכַּלָּה תַּעְדֶּה כֵלֶיהָ׃
    וְיִבְטְחוּ בְךָ יוֹדְעֵי שְׁמֶךָ כִּי לֹא עָזַבְתָּ דֹרְשֶׁיךָ יְהוָה׃
    שִׂמְחוּ בַיהוָה וְגִילוּ צַדִּיקִים וְהַרְנִינוּ כָּל יִשְׁרֵי לֵב׃
    נֵר לְרַגְלִי דְבָרֶךָ וְאוֹר לִנְתִיבָתִי׃
    אוֹדְךָ כִּי עֲנִיתָנִי וַתְּהִי לִי לִישׁוּעָה׃
    
    {\centering
    {\scriptsize    ואומר השיר שהיו הלוים אומרים על הדוכן בכל יום}
    
      \larger ביציאתו מבית המדרש אומר

    }

\begin{tp}{74}
    מוֹדֶה אֲנִי לְפָנֶיךָ יְיָ אֱלֹהַי שֶׁשַּׂמְתָּ חֶלְקִי מִיּוֹשְׁבֵי בֵית הַמֶּדְרָשׁ וְלֹא שַמְתָּ חֶלְקִי מִיּוֹשְׁבֵי קְרָנוֹת
    שֶׁאֲנִי מַשְׁכִּים וְהֵם מַשְׁכִּימִים אֲנִי מַשְׁכִּים לְדִבְרֵי תוֹרָה וְהֵם מַשְׁכִּימִים לִדְבָרֵים בְּטֵלִים 
    אֲנִי עָמֵל וְהֵם עֲמֵלִים אֲנִי עָמֵל וּמְקַבֵּל שָׂכָר וְהֵם עֲמֵלִים וְאֵינָם מְקַבְּלִים שָׂכָר
    אֲנִי רָץ וְהֵם רָצִים אֲנִי רָץ לְחַיֵּי הָעוֹלָם הַבָּא וְהֵם רָצִים לִבְאֵר שַׁחַת\hdot
    שֶׁנֶּאֱמֶר וְאַתָּה אֱלֹהִים תּוֹרִדֵם לִבְאֵר שַׁחַת אַנְשֵׁי דָמִים וּמִרְמָה לֹא יֶחֱצוּ יְמֵיהֶם וַאֲנִי אֶבְטַח בָּךְ׃
\end{tp}


{\centering
  \LARGE תם ונשלם ספר ראשון\\שבחי תודה לאל עליון
}

}

%
%{\centering
%  \LARGE תם ונשלם ספר שלישי\\שבחי תודה לאלהי ישעי
%}

\end{hebrew}
\end{document}
